
\documentclass{article} % For LaTeX2e
\usepackage{iclr2026_conference,times}

\usepackage{xcolor}

% Optional math commands from https://github.com/goodfeli/dlbook_notation.
\input{math_commands.tex}

\usepackage{hyperref}
\usepackage{url}
\usepackage[utf8]{inputenc} % allow utf-8 input
\usepackage[T1]{fontenc}    % use 8-bit T1 fonts
\usepackage{booktabs}       % professional-quality tables
\usepackage{amsfonts}       % blackboard math symbols
\usepackage{nicefrac}       % compact symbols for 1/2, etc.
\usepackage{microtype}      % microtypography
\usepackage{xcolor}         % colors
\usepackage{amsmath}
\usepackage{amssymb}
\usepackage{amsthm}
\usepackage{natbib}
\usepackage{graphicx}
\usepackage{subcaption}
\usepackage{booktabs}   
\usepackage{float}  
\usepackage{placeins}

% amsmath, amssymb, natbib, graphicx, url, algorithm2e

\newtheorem{theorem}{Theorem}[section]
\newtheorem{definition}[theorem]{Definition}
\newtheorem{proposition}[theorem]{Proposition}
\newtheorem{lemma}[theorem]{Lemma}
\newtheorem{assumption}[theorem]{Assumption}
\newtheorem{remark}[theorem]{Remark}
\newtheorem{corollary}[theorem]{Corollary}

\title{Generalization of Gibbs and Langevin Monte Carlo Algorithms in the Interpolation Regime}


% Authors must not appear in the submitted version. They should be hidden
% as long as the \iclrfinalcopy macro remains commented out below.
% Non-anonymous submissions will be rejected without review.

\author{Antiquus S.~Hippocampus, Natalia Cerebro \& Amelie P. Amygdale \thanks{Use footnote for providing further information
about author (webpage, alternative address)---\emph{not} for acknowledging
funding agencies.  Funding acknowledgements go at the end of the paper.} \\
Department of Computer Science\\
Cranberry-Lemon University\\
Pittsburgh, PA 15213, USA \\
\texttt{\{hippo,brain,jen\}@cs.cranberry-lemon.edu} \\
\And
Ji Q. Ren \& Yevgeny LeNet \\
Department of Computational Neuroscience \\
University of the Witwatersrand \\
Joburg, South Africa \\
\texttt{\{robot,net\}@wits.ac.za} \\
\AND
Coauthor \\
Affiliation \\
Address \\
\texttt{email}
}

% The \author macro works with any number of authors. There are two commands
% used to separate the names and addresses of multiple authors: \And and \AND.
%
% Using \And between authors leaves it to \LaTeX{} to determine where to break
% the lines. Using \AND forces a linebreak at that point. So, if \LaTeX{}
% puts 3 of 4 authors names on the first line, and the last on the second
% line, try using \AND instead of \And before the third author name.

% Todonotes is useful during development; simply uncomment the next line
%    and comment out the line below the next line to turn off comments
% \usepackage[disable,textsize=tiny]{todonotes}
\usepackage[textwidth=1.85cm,textsize=tiny]{todonotes}
\newcommand{\massi}[2][]{\todo[color=brown!20,#1]{{\bf MP:} #2}}
\newcommand{\erfan}[2][]{\todo[color=blue!20,#1]{{\bf EM:} #2}}
\newcommand{\andreas}[2][]{\todo[color=green!20,#1]{{\bf AM:} #2}}

\newcommand{\MP}[1]{\textcolor{brown}{#1}}
\newcommand{\EM}[1]{\textcolor{blue}{#1}}
\newcommand{\AM}[1]{\textcolor{teal}{#1}}


\newcommand{\fix}{\marginpar{FIX}}
\newcommand{\new}{\marginpar{NEW}}

% \iclrfinalcopy % Uncomment for camera-ready version, but NOT for submission.
\begin{document}


\maketitle

\begin{abstract}
The paper provides data-dependent bounds on the test error of the Gibbs algorithm in the overparameterized interpolation regime, where low training errors are also obtained for impossible data, such as random labels in classification. The bounds are stable under approximation with Langevin Monte Carlo algorithms. Experiments on the MNIST and CIFAR-10 datasets verify that the bounds yield nontrivial predictions on true labeled data and correctly upper bound the test error for random labels. Our method indicates that generalization in the low-temperature, interpolation regime is already signaled by small training errors in the more classical high temperature regime.
\end{abstract}

\section{Introduction\label{Section introduction}}
Modern learning algorithms can achieve very small training errors on arbitrary data if the underlying hypothesis space is large enough. For reasonable data originating from real-world problems, the chosen hypotheses also tend to have small test errors, a fortunate circumstance, which has given great technological and economic thrust to deep learning.
Unfortunately, the same algorithms also achieve very small training errors for data specifically designed to produce very large test errors, such as random labels in classification. Consequently, the hypothesis space and the training error do not suffice to predict
the test error. The key to generalization must be more deeply buried in the data. \\
While not so disquieting to practitioners, this mystery has troubled theoreticians for many years \citep{zhang2016understanding,zhang2021understanding}, and it seems safe to say that the underlying mechanisms still have not been completely understood.

We are far from solving this riddle in generality, but for the Gibbs posterior we show how nontrivial bounds on the test error can be recovered from the training data. The Gibbs posterior assigns probabilities, which decrease exponentially with the training error of the hypotheses. The exponential decay parameter $\beta $ can be interpreted as an inverse temperature in an analogy to statistical physics. The Gibbs measure is a sufficient idealization to have tractable theoretical properties, but also
the limiting distribution of several concrete stochastic
algorithms, here summarized as Langevin Monte Carlo (LMC), including
Stochastic Gradient Langevin Dynamics (SGLD), \citep{gelfand1991recursive,welling2011bayesian}, a popular modern learning algorithm. 

When $\beta$ is large and the hypothesis space is rich, these algorithms can reproduce the dilemma described above by achieving very small training errors on data designed to have large test errors. 
%If $\beta $ is large, and there are enough hypotheses, these algorithms can reproduce the above-mentioned dilemma by also achieving very small training errors for data designed to have large test errors. 
Our paper addresses this \textit{interpolation regime} of the Gibbs posterior and makes the following three
contributions:
\begin{itemize}
    \item We give high-probability data-dependent bounds on the test error, both for a hypothesis drawn from the Gibbs posterior and for the posterior mean, assuming that we can freely draw samples from the posterior.  These bounds hold for the entire range of temperatures.
    \item We weaken the above assumption by showing that the bounds are stable
    under approximations of the posterior in the total variation and $W_{2}$%
    -Wasserstein metrics. Given sufficient computing resources, this yields bounds for LMC algorithms.  
    \item Using LMC on real-world classification data (MNIST and CIFAR-10), we obtain non-trivial upper bounds on the test error for true labels and
    upper bounds for random labels, despite the algorithms achieving very small training error on the random labels.
%    despite a very small training error on random labels.
\end{itemize}
Our method is based on a combination of the PAC-Bayesian bounds \citep{mcallester1999pac,alquier2021user,rivasplata2020pac} with
an integral representation of the log-partition function. This makes it possible to
bound the logarithm of the density of the Gibbs posterior at a given temperature in terms of empirical averages at higher temperatures. A qualitative conclusion is that generalization in the under-regularized low-temperature regime ($\beta >n$) is already indicated by small training errors in the over-regularized high-temperature regime ($\beta <n$), where $n$ is the number of training examples.

\subsection{Related literature}

Many papers address the generalization of the Gibbs algorithm and Langevin Monte Carlo, with special focus on SGLD, which is the most popular
algorithm. Most similar to this work is \cite{raginsky2017non}, which bounds
the distance to the Gibbs posterior and then its generalization
error. Their bound, however, applies only to the high temperature regime $%
\beta <n$.

Several works concentrate on the optimization path of SGLD. \cite%
{mou2018generalization} gives both stability and PAC-Bayesian bounds. \cite%
{pensia2018generalization} applies the information theoretic generalization
bounds of \cite{xu2017information}. These ideas are further developed by 
\cite{negrea2019information}, where random subsets of the training data
are used to define data-dependent priors. \cite{farghly2021time} gives time-independent bounds for SGLD, which are further improved by \cite%
{futami2024time}. Most of the bounds in the above papers are in expectation.
The very recent paper of \cite{harel2025temperature} gives a very elegant
argument for Markov chain algorithms based on the second law of
thermodynamics. If the invariant distribution is the Gibbs posterior, the
bound along the entire optimization path is of order $\sqrt{\beta /n}$ but
improvable to $\beta /n$.

Some papers give similar bounds for the Gibbs posterior, roughly of the form $\beta /n$ or $\sqrt{\beta /n}$ \citep{raginsky2017non,dziugaite2018data,kuzborskij2019distribution,rivasplata2020pac} or \cite{maurer2024generalization} and \cite%
{harel2025temperature}). These bounds hold equally for random labels and are therefore vacuous for overparametrized hypothesis spaces in the low
temperature regime $\beta >n$. To our knowledge, ours is the only bound for the Gibbs posterior, which is valid
in this regime.

Other bounds have been developed for specific algorithms  designed to optimize them. The milestone paper by \cite{dziugaite2017computing} is the most prominent example, and  \citep{dziugaite2018data}
and \citep{perez2021tighter} are also in this category. Our bounds by contrast apply to the Gibbs posterior and LMC in
their standard forms.

\section{Preliminaries}\label{Section preliminaries}

The relative entropy of two Bernoulli variables with expectations $p$ and $q$
is denoted 
\begin{equation}
\kappa \left( p,q\right) =p\ln \frac{p}{q}+\left( 1-p\right) \ln \frac{1-p}{%
1-q}.  \label{define little kl}
\end{equation}%
 We also define the function $%
\kappa ^{-1}:\left[ 0,1\right] \times \left[ 0,\infty \right) \rightarrow %
\left[ 0,1\right] $ by%
\begin{equation*}
\kappa ^{-1}\left( p,t\right) =\inf \left\{ q:q\geq p,\kappa \left(
p,q\right) \geq t\right\} .
\end{equation*}

Throughout the following $\left( \mathcal{X},\Sigma \right) $ is a
measurable space of \textit{data} with probability measure $\mu $. The iid
random vector $\mathbf{x}\sim \mu ^{n}$ is the training sample. \\ We let $\left( \mathcal{H},\Omega \right) $ be a measurable space of \textit{%
hypotheses}, and let $\ell :\mathcal{%
H\times X}\rightarrow \left[ 0,\infty \right)$ be a prescribed 
%there is a fixed measurable 
loss function. Members of $\mathcal{H}$
are denoted $h$ or $g$. We write $L\left( h\right) :=\mathbb{E}_{x\sim \mu }%
\left[ \ell \left( h,x\right) \right] $ and $\hat{L}\left( h,\mathbf{x}%
\right) :=\left( 1/n\right) \sum_{i}\ell \left( h,x_{i}\right) $
respectively for the true (expected) and empirical error of hypothesis $h\in 
\mathcal{H}$. The set of probability measures on $\left( \mathcal{H},\Omega \right) $ is
denoted $\mathcal{P}\left( \mathcal{H}\right) $. \\ A stochastic algorithm is a
function $\nu :\mathcal{X}^{n}\rightarrow \mathcal{P}\left( \mathcal{H}%
\right) $, which assigns to a training sample $\mathbf{x}$ a probability
measure $\nu \left( \mathbf{x}\right) \in \mathcal{P}\left( \mathcal{H}%
\right) $. The $KL$-divergence between two probability measures is the
function $KL:(\rho,\nu)\in \mathcal{P}( \mathcal{H})\times \mathcal{P}(\mathcal{H})\mapsto \mathbb{E}
_{h\sim \rho}\big[ \ln \frac{d\rho }{d\nu }\big] $ if $\rho $ absolutely
continuous w.r.t. $\nu $, otherwise the value is $\infty $. The total
variation distance is defined as $d_{TV}:( \rho,\nu) \in 
\mathcal{P}(\mathcal{H}) \times \mathcal{P}(\mathcal{H})\mapsto \sup_{A\in \Omega }\vert \rho(A) -\nu(A)\vert$. The $W_{p}$-Wasserstein distance is $
W_{p}(\rho,\nu) =( \inf_{W}\mathbb{E}_{(x,y)
\sim W}[\Vert x-y\Vert^{p}]) ^{1/p}$ with the
infimum over all probability measures on $\mathcal{P}( \mathcal{H}) \times \mathcal{P}( \mathcal{H}) $ with $\rho $ and $\nu $
as marginals.

There is an a-priori reference measure $\pi \in \mathcal{P}\left( \mathcal{H}%
\right) $, called the \textit{prior}. With a fixed prior, the Gibbs algorithm
at inverse temperature $\beta >0$ is the stochastic algorithm $G_{\beta }:%
\mathbf{x}\in \mathcal{X}^{n}\mapsto G_{\beta }\left( \mathbf{x}\right) \in 
\mathcal{P}\left( \mathcal{H}\right) $ defined by%
\begin{equation*}
G_{\beta }\left( \mathbf{x}\right) \left( A\right) =\frac{1}{Z_{\beta
}\left( \mathbf{x}\right) }\int_{A}e^{-\beta \hat{L}\left( h,\mathbf{x}%
\right) }d\pi \left( h\right) \text{ for }A\in \Omega \text{.}
\end{equation*}%
$G_{\beta }\left( \mathbf{x}\right) $ is called the \textit{Gibbs-posterior, 
}the normalizing factor 
\begin{equation*}
Z_{\beta }\left( \mathbf{x}\right) :=\int_{\mathcal{H}}e^{-\beta \hat{L}%
\left( h,\mathbf{x}\right) }d\pi \left( h\right)
\end{equation*}%
is called the \textit{partition function}. The motivation for the Gibbs posterior is that it puts larger weights on hypotheses with smaller empirical error.\newline 
Given a stochastic algorithm $\nu $ we define a probability measure $\rho
_{\nu }$ on $\mathcal{H}\times \mathcal{X}^{n}$ by 
\begin{equation}
\rho _{\nu }\left( A\right) =\mathbb{E}_{\mathbf{x}\sim \mu ^{n}}\mathbb{E}%
_{h\sim \nu \left( \mathbf{x}\right) }\left[ 1_{A}\left( h,\mathbf{x}\right) %
\right] \text{ for }A\in \Omega \otimes \Sigma ^{\otimes n}.
\label{Define rho}
\end{equation}%
Then, $\mathbb{E}_{\left( h,\mathbf{x}\right) \sim \rho _{\nu }}\left[ \phi
\left( h,\mathbf{x}\right) \right] =\mathbb{E}_{\mathbf{x}}\mathbb{E}_{h\sim
\nu \left( \mathbf{x}\right) }\left[ \phi \left( h,\mathbf{x}\right) \right] 
$ for measurable $\phi :\mathcal{H}\times \mathcal{X}^{n}\rightarrow \mathbb{%
R}$. To draw the pair $\left( h,\mathbf{x}\right) $ from $\rho _{\nu }$ we
first draw the training sample $\mathbf{x}$, and then sample $h$ from $\nu
\left( \mathbf{x}\right) $. The main objective in learning is that the risk $%
\mathbb{E}_{x\sim \mu }\left[ f\left( h,x\right) \right] $ is small with
high probability in $\left( h,\mathbf{x}\right) \sim \rho _{\nu }$, where $f$
is some application-dependent loss function, possibly different from $\ell $. In the sequel we will give corresponding guarantees.
\section{Bounds for the Gibbs posterior}

In this section, we make the idealized assumption that we are free to sample
from the Gibbs posterior at any finite $\beta \geq 0$.

\subsection{An integral representation of the free energy\label{Section integral representation}}

\begin{lemma}
\label{Lemma Derivatives}Let $0=\beta _{0}<\beta _{1}<\cdots<\beta _{K}=\beta $%
. Then%
\begin{equation*}
-\ln Z_{\beta }(\mathbf{x}) =\int_{0}^{\beta }\mathbb{E}_{h\sim
G_{\gamma }(\mathbf{x})}\big[ \hat{L}( h,\mathbf{x})\big] d\gamma \leq \sum_{k=1}^{K}(\beta_{k}-\beta_{k-1})\mathbb{E}_{g\sim G_{\beta _{k-1}}(\mathbf{x})}\big[ \hat{L}( g,\mathbf{x}) \big].
\end{equation*}
\end{lemma}
\begin{proof}
Let $A(\beta) =-\ln Z_{\beta}(\mathbf{x})$. One
verifies the identities
\begin{eqnarray*}
A(0) &=&0, \\
A^{\prime}(\beta) &=&\frac{1}{Z_{\beta,\pi}(\mathbf{x})}\int_{\mathcal{H}}\hat{L}(h,\mathbf{x})e^{-\beta\hat{L}(h,\mathbf{x}) }d\pi(h)=\mathbb{E}_{h\sim G_{\beta}(\mathbf{x}}\big[\hat{L}(h,\mathbf{x})
\big],\\
A^{\prime \prime}(\beta) &=&-\Big( \mathbb{E}_{h\sim G_{\beta
}(\mathbf{x})}\big[\hat{L}(h,\mathbf{x})^{2}
\big] -\big(\mathbb{E}_{h\sim G_{\beta}(\mathbf{x})}\big[\hat{L}(h,\mathbf{x}) \big] \big)^{2}\Big) \leq 0\text{.}
\end{eqnarray*}
The equality in the lemma then follows from the first two identities above and the fundamental theorem of calculus, and the inequality follows from the last identity, which shows that $\mathbb{E}_{g\sim G_{\beta _{k-1},\pi
}(\mathbf{x})}\big[\hat{L}( g,\mathbf{x})\big]$
is non-increasing in $\beta$.
\end{proof}
In statistical physics there is a formal analogy, where the function $h\mapsto \hat{L}(h,\mathbf{x})$ is the ($\mathbf{x}$-dependent) energy of the system in the state $h$, and $\beta$ is the inverse
temperature. The Gibbs posterior then becomes the "canonical ensemble" \citep{gibbs1902elementary},
describing the probability of states in equilibrium with a heat bath at temperature $\beta ^{-1}$. The function $\beta
\mapsto A(\beta)$ plays an important role: $\beta ^{-1}A(\beta) $ is the Helmholz free energy, $A^{\prime}(\beta)=\mathbb{E}_{h\sim G_{\beta}(\mathbf{x})}\big[\hat{L}(h,\mathbf{x})\big]$ is the thermal average of the energy, $-\beta A^{\prime }(\beta)+A(\beta) $ is the entropy and $-A^{\prime \prime}(\beta )$ is proportional to the heat capacity at temperature $\beta ^{-1}$ \citep[see e.g.][]{huang2008statistical}.

For $h\in \mathcal{H}$, $\mathbf{x\in }\mathcal{X}^{n}$ and an increasing sequence $\bm{\beta}=\left( \beta _{1}<\cdots<\beta _{K}\right) 
$ of positive numbers, we denote \begin{equation}
\Gamma(h,\mathbf{x},\bm{\beta})=-\beta_{K}\hat{L}(h,\mathbf{x})+\sum_{k=1}^{K}(\beta_{k}-\beta_{k-1})\mathbb{E}_{g\sim G_{\beta_{k-1}}(\mathbf{x})}\big[\hat{L}(g,\mathbf{x})\big].
\label{eq:gamma}
\end{equation}%
So Lemma \ref{Lemma Derivatives} states that, when $\beta_{0}=0$, then,  
\begin{equation}
-\beta _{K}\hat{L}\left( h,\mathbf{x}\right) -\ln Z_{\beta _{K}}\left( 
\mathbf{x}\right) \leq \Gamma \left( h,\mathbf{x},\bm{\beta}\right) .
\label{density bound}
\end{equation}%
Note that $\Gamma \left( h,\mathbf{x},\bm{\beta}\right) $ depends 
\textit{only} on the training data $\mathbf{x}$, the sequence $\bm{\beta 
}$ and the hypothesis $h$.

\subsection{Bounds\label{Section bounds}}
The function $F$ in the following is a placeholder for a random variable related to the generalization gap, which we would like to bound with high probability.
\begin{theorem}
\label{Theorem Main}Let $F:\mathcal{H}\times \mathcal{X}^{n}\rightarrow 
\mathbb{R}$ be some measurable function, $\beta >0$ and $\bm{\beta}=\left( \beta _{1}<\cdots<\beta _{K}\right) $ as above with $\beta_{0}=0$ and $\beta _{K}=\beta $. Then,

(i) for $\delta >0$ with probability at least $1-\delta $ in $\mathbf{x}\sim
\mu ^{n}$ and $h\sim G_{\beta }\left( \mathbf{x}\right) $ 
\begin{equation*}
F\left( h,\mathbf{x}\right) \leq \Gamma \left( h,\mathbf{x},\bm{\beta}%
\right) +\ln \mathbb{E}_{\mathbf{x}}\mathbb{E}_{g\sim \pi }\left[ e^{F\left(
g,\mathbf{x}\right) }\right] +\ln \left( 1/\delta \right),
\end{equation*}

(ii) for $\delta >0$ with probability at least $1-\delta $ in $\mathbf{x}%
\sim \mu ^{n}$ 
\begin{equation*}
\mathbb{E}_{h\sim G_{\beta }\left( \mathbf{x}\right) }\left[ F\left( h,%
\mathbf{x}\right) \right] \leq \mathbb{E}_{h\sim G_{\beta }\left( \mathbf{x}%
\right) }\left[ \Gamma \left( h,\mathbf{x},\bm{\beta}\right) \right]
+\ln \mathbb{E}_{\mathbf{x}}\mathbb{E}_{g\sim \pi }\left[ e^{F\left( g,%
\mathbf{x}\right) }\right] +\ln \left( 1/\delta \right).
\end{equation*}
\end{theorem}

\begin{proof}
By Markov's inequality, for any real random variable $Y$ 
\begin{equation*}
\Pr \left\{ Y>\ln \mathbb{E}\left[ e^{Y}\right] +\ln \left( 1/\delta \right)
\right\} =\Pr \left\{ e^{Y}>\mathbb{E}\left[ e^{Y}\right] /\delta \right\}
\leq \delta .
\end{equation*}%
To prove (i), we apply this to the random variable $Y=F\left( h,\mathbf{x}%
\right) +\beta \hat{L}\left( h,\mathbf{x}\right) +\ln Z_{\beta }\left( 
\mathbf{x}\right) $ on the probability space $\left( \mathcal{H\times X}%
^{n},\Omega \otimes \Sigma ^{\otimes n},\rho_{G_\beta} \right) $ as defined in (\ref%
{Define rho}). Together with the definition of the Gibbs posterior, this
gives, with probability at least $1-\delta $ in $\left( h,\mathbf{x}\right)
\sim \rho_{G_\beta} $ (equivalent to saying $\mathbf{x}\sim \mu ^{n}$ and $h\sim
G_{\beta }\left( \mathbf{x}\right) $),%
\begin{align*}
& F\left( h,\mathbf{x}\right) +\beta \hat{L}\left( h,\mathbf{x}\right) +\ln Z_{\beta }\left( \mathbf{x}\right) \\ & \leq \ln \mathbb{E}_{\mathbf{x}}\mathbb{E}_{g\sim G_{\beta }\left( \mathbf{x}\right) }\left[ e^{F\left( g,\mathbf{x}\right) +\beta \hat{L}\left( g,\mathbf{x}\right) +\ln Z_{\beta }\left( \mathbf{x}\right) }\right] +\ln\left( 1/\delta \right) \\& =\ln \mathbb{E}_{\mathbf{x}}\mathbb{E}_{g\sim \pi }\left[ e^{F\left( g,\mathbf{x}\right) +\beta \hat{L}\left( g,\mathbf{x}\right) +\ln Z_{\beta}\left( \mathbf{x}\right) -\beta \hat{L}\left( g,\mathbf{x}\right) -\ln Z_{\beta }\left( \mathbf{x}\right) }\right] +\ln \left( 1/\delta \right) \\ & =\ln \mathbb{E}_{\mathbf{x}}\mathbb{E}_{g\sim \pi }\left[ e^{F\left( g, \mathbf{x}\right) }\right] +\ln \left( 1/\delta \right).
\end{align*}
Subtract $\beta \hat{L}\left( h,\mathbf{x}\right) +\ln Z_{\beta }\left( 
\mathbf{x}\right) $ and use (\ref{density bound}). For (ii)\ apply Markov's
inequality to $\mathbb{E}_{h\sim G_{\beta }\left( \mathbf{x}\right) }\left[
F\left( h,\mathbf{x}\right) +\beta \hat{L}\left( h,\mathbf{x}\right) +\ln
Z_{\beta }\left( \mathbf{x}\right) \right] $ instead. By Jensen's inequality 
\begin{equation*}
e^{\mathbb{E}_{h\sim G_{\beta }\left( \mathbf{x}\right)}\left[ F\left(h,\mathbf{x}\right) +\beta \hat{L}\left( h,\mathbf{x}\right) +\ln Z_{\beta
}\left(\mathbf{x}\right)\right] }\leq \mathbb{E}_{g\sim G_{\beta}\left( 
\mathbf{x}\right) }\left[e^{F\left( g,\mathbf{x}\right) +\beta \hat{L}
\left( g,\mathbf{x}\right)+\ln Z_{\beta }\left( \mathbf{x}\right)}\right]
\end{equation*}%
and proceed as above using (\ref{density bound}).
\end{proof}

Up to the application of (\ref{density bound}), the above proof of (i) just
gives the single-draw version of the PAC-Bayesian bound as in \cite{rivasplata2020pac} applied to the Gibbs posterior, while (ii) is the
standard PAC-Bayesian bound applied to the Gibbs posterior until (\ref{density bound}) is invoked.

\subsection{Loss functions and secondary loss functions\label{Section
secondary loss}}

To apply Theorem \ref{Theorem Main}, we need to control the
exponential moment $\mathbb{E}_{\mathbf{x}}\mathbb{E}_{g\sim \pi }\left[
e^{F\left( g,\mathbf{x}\right) }\right] $, but otherwise we have free choice
of the function $F$. This gives the method some flexibility. If $\exp \left(
F\right) $ has sufficient integrability properties, we can exchange the two
expectations and often there is a bound on $\mathbb{E}_{\mathbf{x}}\left[
e^{F\left( g,\mathbf{x}\right) }\right] $ uniform in $g$, which then carries
over to $\mathbb{E}_{g\sim \pi }\mathbb{E}_{\mathbf{x}}\left[ e^{F\left( g,%
\mathbf{x}\right) }\right] $, since $\pi $ is a probability measure. In this way bounds for sub-Gaussian or sub-exponential losses can be
obtained, but also for U-statistics or even non-iid data, sampled from the
trajectories of time-homogeneous, ergodic Markov chains. In Section \ref%
{Section Subgaussian} in the appendix, we derive a bound for sub-Gaussian
losses from Theorem \ref{Theorem Main}; other examples are planned for a
longer version of the paper.

The function $F$ may be defined in terms of other, application-dependent
loss functions, which are different from the loss $\ell $, which defines the
Gibbs posterior and the functional $\Gamma $. To illustrate this point, let $f:\mathcal{H}\times\mathcal{X}\rightarrow [0,1] $ be measurable
and set $F(h,\mathbf{x}) =n~\kappa \left(\frac{1}{n}\sum_{i}f(h,x_{i}) ,\mathbb{E}_{x}\left[f(h,x) \right] \right) $, with $\kappa $ the relative entropy as in (\ref{define
little kl}). Then, Theorem 1 of \cite{maurer2004note} gives $\mathbb{E}_{\mathbf{x}}\left[e^{F(h,\mathbf{x})}\right]\leq 2\sqrt{n}$ for $n\geq 8$. Substitution in Theorem \ref{Theorem Main} and division by $n$ then give the following corollary.

\begin{corollary}
\label{Corollary kl}Let $f:\mathcal{H}\times\mathcal{X}\rightarrow [0,1]$ be measurable, $\delta>0$ and $n\geq 8$. Then, with probability
at least $1-\delta$ in $\mathbf{x}\sim \mu ^{n}$ and $h\sim G_{\beta
}(\mathbf{x})$
\begin{equation*}
\kappa \left(\frac{1}{n}\sum_{i}f(h,x_{i}),\mathbb{E}_{x}[f(h,x)]\right) \leq \frac{1}{n}\left(\Gamma(h,\mathbf{x},\bm{\beta})+\ln \left(\frac{2\sqrt{n}}{\delta}\right) \right)
\end{equation*}
and with probability at least $1-\delta $ in $\mathbf{x}\sim \mu ^{n}$%
\begin{equation*}
\kappa \left(\frac{1}{n}\sum_{i} \mathbb{E}_{h\sim G_{\beta }\left( \mathbf{%
x}\right) }\left[ f\left( h,x_{i}\right) \right] ,\mathbb{E}_{h\sim G_{\beta
}\left( \mathbf{x}\right) }\mathbb{E}_{x}\left[ f\left( h,x\right) \right]
\right) \leq \frac{1}{n}\left( \Gamma \left( h,\mathbf{x},\bm{\beta}%
\right) +\ln \left( \frac{2\sqrt{n}}{\delta }\right) \right) .
\end{equation*}
\end{corollary}

For the second part, we used the joint convexity of $\kappa $. Under the
conditions of this corollary, the second inequality becomes%
\begin{equation*}
\mathbb{E}_{h\sim G_{\beta }\left( \mathbf{x}\right) }\mathbb{E}_{x}\left[
f\left( h,x\right) \right] \leq \kappa ^{-1}\left( \frac{1}{n}\sum_{i}%
\mathbb{E}_{h\sim G_{\beta }\left( \mathbf{x}\right) }\left[ f\left(
h,x_{i}\right) \right] ,\frac{1}{n}\left( \Gamma \left( h,\mathbf{x},\mathbf{%
\beta }\right) +\ln \left( \frac{2\sqrt{n}}{\delta }\right) \right) \right)
\end{equation*}%
with an analogous version for the single-draw case. This is how we compute bounds in our experiments. For an illustration, please refer to \ref{sec:illustration}. 

Here $f$ plays the role of a secondary loss function, typically different from the loss function $\ell$, which defines the Gibbs posterior.
In applications one would define the Gibbs posterior in terms of a differentiable loss function $\ell $ and approximate it with a suitable Monte Carlo method. For classification, however, one is interested in bounding the 0-1 loss obtained by comparing $\ell $ to some threshold $\theta$. Then (momentarily changing notation by replacing $x\in \mathcal{X}$
by $\left( x,y\right) $, where $y$ is the label corresponding to $x$) one
can define $f\left( h,\left( x,y\right) \right) =1_{\left( -\infty ,\theta
\right) }\left( y\ell \left( h,x\right) \right) $ to obtain a bound on the
expected 0-1 loss. 

\section{Bounds for Langevin Monte Carlo\label{Langevin Theory}}

For this section, we assume $\mathcal{H}=\mathbb{R}^{d}$ and an isotropic Gaussian prior $\pi$ of width $\sigma$. We condition on the
training data $\mathbf{x}$, reference to which we omit. The Gibbs posterior
is an idealization, from which it is impossible to sample directly.
Nevertheless a number of works \citep{raginsky2017non,
dalalyan2017user,brosse2018promises,vempala2019rapid,dwivedi2019log,nemeth2021stochastic,balasubramanian2022towards} discuss algorithms (SGLD, ULA, MALA, etc, here
summarized as Langevin Monte Carlo (LMC)), capable of approximating a
probability measure $\nu$ on $\mathbb{R}^{d}$ of the form $\nu \propto \exp
\left( -V\right) $ or some nearby limiting distribution. In the following, we discuss one of these algorithms. 
\subsection{ULA\label{Section ULA}}
We focus on the results of \cite{vempala2019rapid}, which do not
require convexity of $V$ and instead assume that the measure $\nu $
satisfies a log-Sobolev inequality (LSI) in the sense that for all smooth $f:%
\mathbb{R}^{d}\rightarrow \mathbb{R}$%
\begin{equation}
\mathbb{E}_{h\sim \nu }\left[ f^{2}\left( h\right) \ln f^{2}\left( h\right) %
\right] -\mathbb{E}_{h\sim \nu }\left[ f^{2}\left( h\right) \right] \ln 
\mathbb{E}_{h\sim \nu }\left[ f^{2}\left( h\right) \right] \leq \frac{2}{%
\alpha }\mathbb{E}_{h\sim \nu }\left[ \left\Vert \left( \nabla f\right)
\left( h\right) \right\Vert ^{2}\right]  \label{log sobolev}
\end{equation}%
for some $\alpha >0$. An LSI is satisfied when $V$ is strongly convex, but,
importantly, also for measures which are bounded perturbations of measures
satisfying an LSI (\cite{holley1986logarithmic}). \cite{vempala2019rapid}
give further examples and a list of references for measures, which are not
log-concave and satisfy an LSI. \cite{raginsky2017non} show, that
under dissipativity conditions of the loss the Gibbs posterior $G_{\beta
}\left( \mathbf{x}\right) $ satisfies an LSI with constant independent of $%
\mathbf{x}$.

Consider the iterative algorithm 
\begin{equation}
h_{t+1}=h_{t}-\epsilon \nabla V\left( h_{t}\right) +\sqrt{2\epsilon }\xi
_{t},  \label{LMC}
\end{equation}%
where $\epsilon $ is a step size, the $\xi _{t}\sim \mathcal{N}\left(
0,I\right) $ are independent Gaussian vectors and $h_{0}$ is drawn from some
initial distribution $\nu _{0}$. Some authors call this algorithm simply
LMC, for Langevin Monte Carlo. We call it ULA, alongside \cite{durmus2017nonasymptotic},\cite{dwivedi2019log} and \cite{vempala2019rapid}, for Un-adjusted Langevin Algorithm, because it misses the Metropolis-type accept-reject step, which would guarantee that the invariant distribution is indeed the Gibbs posterior. A popular variant of ULA is Stochastic Gradient Langevin Dynamics (SGLD) \citep{welling2011bayesian,raginsky2017non} where the gradient is replaced by an
unbiased estimate, typically realized with random minibatches. Here, we restrict ourselves to ULA with a constant step size, because it has
the least number of parameters to adjust, but in experiments we also use the computationally more efficient SGLD.

As $\epsilon \,{\rightarrow}\, 0$, ULA recovers the Continuous Langevin Dynamics
(CLD) given by the stochastic differential equation
\begin{equation*}
dh_{t}=- \nabla V\left( h_{t}\right) dt+\sqrt{2}dB_{t},
\end{equation*}%
where $B_{t}$ is centered standard Brownian motion in $\mathbb{R}^{d}$. CLD
converges exponentially to the Gibbs posterior \citep{chiang1987diffusion}.
For $\epsilon >0$, the distribution $\nu _{\epsilon ,t}$ of ULA converges as $
t\rightarrow \infty $ to a biased limiting distribution $\nu _{\epsilon }$
which is generally different from $\nu $, but expected to be closer to $\nu $
as $\epsilon $ becomes smaller. \cite{vempala2019rapid} use the LSI
assumption and coupling to control the difference between CLD and ULA along
their path and prove the following result.

%\massi{refer to Vempala and W in the theorem? e.g. as we do in Lemma A.2}
\begin{theorem}\label{Theorem Vempala}
Assume that $\nu $ satisfies the log-Sobolev
inequality (\ref{log sobolev}) with $\alpha >0$, that the Hessian of $V$
satisfies $-LI\preceq \nabla ^{2}V\left( h\right) \preceq LI$ for all $h$
and some $L<\infty $, and that $0<\epsilon \leq \alpha /\left( 4L^{2}\right) 
$. Then, for $t\geq 0$%
\begin{equation*}
KL\left( \nu ,\nu _{\epsilon ,t}\right) \leq e^{-\alpha \epsilon t}KL\left(
\nu ,\nu _{0}\right) +\frac{8\epsilon dL^{2}}{\alpha }.
\end{equation*}
\end{theorem}

The first exponential term is due to the mismatch of the initial
distribution and $\nu$. The second term bounds the divergence between the
limiting distribution $\nu _{\epsilon }\,\ $and $\nu $. Similar results
exist under different conditions on the potential $V$. \cite{cheng2018sharp}
for example require $V$ to be strongly convex outside of a ball instead of
the log-Sobolev inequality and gives bounds in terms of the $W_{1}$%
-Wasserstein metric. \cite{raginsky2017non} give bounds for $W_{2}$ under dissipativity assumptions. The next corollary adapts Theorem \ref{Theorem Vempala}
to the situation studied in this paper.

%\massi{is the Gaussian prior below a specific choice for the $\pi$ used earlier in the paper? I found this confusing. Is this prior introduced to satisfy LSI? Also in general in this section there are several distributions at play: $G_\beta$, $\nu_0$, $\nu_\epsilon$, $\nu_{\epsilon,t}$. A fig showing how they  relate would help the reader. Probably no time for this but let's discuss. Also eq. (6) is a bit out of the blue}
\begin{corollary}
\label{Corollary approximations}For $\beta >0$ consider the Gibbs posterior $%
G_{\beta }$ corresponding to $\hat{L}\left( h\right) $, with centered
Gaussian prior of width $\sigma $. Assume that it satisfies the log-Sobolev
inequality (\ref{log sobolev}) with $\alpha >0$, that the Hessian of $\hat{L}
$ satisfies $-RI\preceq \nabla ^{2}\hat{L}\left( h\right) \preceq RI$ for
all $h$ and some $R<\infty $, and that $0<\eta \leq \alpha /\left( 4\left(
\beta R+\frac{1}{\sigma ^{2}}\right) ^{2}\right) $. Consider the algorithm
\begin{equation}
h_{t+1}=h_{t}-\eta \nabla _{h}\hat{L}\left( h_{t}\right) -\frac{\eta h_{t}}{%
\beta \sigma ^{2}}+\sqrt{\frac{2\eta }{\beta }}\xi _{t},  \label{Our LMC}
\end{equation}%
where $h_{0}\sim \nu _{0}$ and the $\xi _{t}\sim \mathcal{N}\left(
0,I\right) $ are independent Gaussian random variables. Let $D\left( \beta
\right) =KL\left( G_{\beta },\nu _{0}\right) $ and let $\nu _{\beta ,\eta ,t}
$ be the distribution of $h_{t}$ after $t$ steps. Then,

(i) $KL\left( G_{\beta },\nu _{\beta ,\eta ,t}\right) \leq e^{-\alpha \eta
t/\beta }D\left( \beta \right) +\frac{8\eta d}{\beta \alpha }\left( \beta R+%
\frac{1}{\sigma ^{2}}\right) ^{2}.$

(ii) $W_{2}\left( G_{\beta },\nu _{\beta ,\eta ,t}\right) \leq \frac{2}{%
\alpha }e^{-\alpha \eta t/\beta }D\left( \beta \right) +\frac{16\eta d}{%
\beta \alpha ^{2}}\left( \beta R+\frac{1}{\sigma ^{2}}\right) ^{2}.$

(iii) $d_{TV}\left( G_{\beta },\nu _{\beta ,\eta ,t}\right) \leq e^{-\alpha
\eta t/\left( 2\beta \right) }\sqrt{D\left( \beta \right) }+2\sqrt{\frac{%
\eta d}{\beta \alpha }}\left( \beta R+\frac{1}{\sigma ^{2}}\right) .$
\end{corollary}

\begin{proof}
We make the identifications $V\left( h\right) =\beta \hat{L}\left( h\right)
+\left\Vert h\right\Vert ^{2}/\left( 2\sigma ^{2}\right) $, $\epsilon =\eta
/\beta $ and $L=\beta R+\frac{1}{\sigma ^{2}}$. Then, $\nu =G_{\beta }$ with
Gaussian prior of width $\sigma $, ULA becomes (\ref{Our LMC}) and (i)
follows directly from Theorem \ref{Theorem Vempala}, (ii) follows from \cite%
{otto2000generalization} and the LSI assumption, and (iii) follows from
Pinsker's inequality (see e.g. \cite{Boucheron13}).
\end{proof}
%\massi{add a short discussion on the result and how we are going to use it?}

\subsection{Stability of the bounds\label{Section stability}}
We now show the stability of our bounds for approximation in total variation and 
$W_{2}$-Wasserstein metrics, under boundedness or Lipschitz conditions. Together with Corollary %
\ref{Corollary approximations} this implies bounds for the algorithm defined
in (\ref{Our LMC}).

We assume that there is a target approximation $\nu _{\beta }\left( \mathbf{x%
}\right) $ of $G_{\beta }\left( \mathbf{x}\right) $, for which we want to
compute a high probability bound, either for the single draw version on $%
F\left( h,\mathbf{x}\right) $ as $\mathbf{x}\sim \mu ^{n}$ and $h\sim \nu
_{\beta }\left( \mathbf{x}\right) $, or, for the classical PAC-Bayesian
version, on $\mathbb{E}_{h\sim \nu _{\beta }\left( \mathbf{x}\right) }\left[
F\left( h,\mathbf{x}\right) \right] $ as $\mathbf{x}\sim \mu ^{n}$. It is
not surprising that the single-draw bound will require a much closer
approximation of the Gibbs posterior.

Since the bounding functional $\Gamma \left( h,\mathbf{x},\bm{\beta}%
\right) $ depends on the Gibbs posteriors $G_{\beta _{k}}\left( \mathbf{x}%
\right) $ for $k\in\{ 1,\dots,K-1\}$, we require corresponding approximations $\nu
_{\beta _{k}}$ of $G_{\beta _{k}}\left( \mathbf{x}\right) $ to compute the
bound. To streamline notation, we define 
\begin{equation*}
\Gamma _{\nu }\left( h,\mathbf{x},\bm{\beta}\right) =-\beta \hat{L}%
\left( h,\mathbf{x}\right) +\sum_{k=1}^{K}\left( \beta _{k}-\beta
_{k-1}\right) \mathbb{E}_{g\sim \nu _{\beta _{k-1}}\left( \mathbf{x}\right) }%
\big[ \hat{L}(g,\mathbf{x}) \big]
\end{equation*}%
for $0=\beta _{0}<\beta _{1}<\cdots<\beta _{K}=\beta $ and $\left( \nu
_{0}\left( \mathbf{x}\right) ,\nu _{\beta _{1}}\left( \mathbf{x}\right)
,\cdots,\nu _{\beta _{k-1}}\left( \mathbf{x}\right) \right) \in \mathcal{P}%
\left( \mathcal{H}\right) ^{K}$. Like $\Gamma \left( h,\mathbf{x},\mathbf{%
\beta }\right)$, the functional $\Gamma _{\nu }\left( h,\mathbf{x},\mathbf{%
\beta }\right) $ depends on the training data, but it can also be computed
by repeated execution of the algorithm (\ref{Our LMC}). The next theorem
states the obtained bound in terms of the approximation errors in total
variation.

\begin{theorem}
\label{Theorem BV approximations}Suppose that $\mathcal{H}=\mathbb{R}^{d}$
and that there are numbers $m,M<\infty $ such that for every $\mathbf{x}$ in 
$\mathcal{X}^{n}$ and $h\in \mathcal{H}$ we have $\left\vert \ell \left( h,%
\mathbf{x}\right) \right\vert \leq m$ and $\left\vert F\left( h,\mathbf{x}%
\right) \right\vert \leq M$. Let $0=\beta _{0}<\beta _{1}<\dots<\beta
_{K}=\beta $ and $\nu _{\beta _{k}}\left( \mathbf{x}\right) \in \mathcal{P}%
\left( \mathcal{H}\right) $ be such that $d_{TV}\left( \nu _{\beta
_{k}}\left( \mathbf{x}\right) ,G_{\beta_{k} }\left( \mathbf{x}\right) \right)
=\epsilon _{\beta _{k}}$. Then,

(i) with probability at least $1-\delta $ as $\mathbf{x}\sim \mu ^{n}$ and $
h\sim \nu _{\beta }\left( \mathbf{x}\right) $%
\begin{equation*}
F\left( h,\mathbf{x}\right) \leq \Gamma _{\nu }\left( h,\mathbf{x},\mathbf{%
\beta }\right) +\ln \mathbb{E}_{\mathbf{x}}\mathbb{E}_{h\sim \pi }\left[
e^{F\left( h,\mathbf{x}\right) }\right] +\ln \frac{1}{\delta } +\ln \left( 2e^{M+\beta m}\epsilon _{\beta }\right) +\sum_{k=1}^{K}\left(
\beta_{k}-\beta _{k-1}\right) m\epsilon _{\beta _{k-1}}.
\end{equation*}

(ii) with probability at least $1-\delta $ as $\mathbf{x}\sim \mu ^{n}$ 
\begin{eqnarray*}
\mathbb{E}_{h\sim \nu _{\beta }\left( \mathbf{x}\right) }\left[ F\left( h,%
\mathbf{x}\right) \right] &\leq &\mathbb{E}_{h\sim \nu _{\beta }\left( 
\mathbf{x}\right) }\left[ \Gamma _{\nu }\left( h,\mathbf{x},\bm{\beta}%
\right) \right] +\ln \mathbb{E}_{\mathbf{x}}\mathbb{E}_{h\sim \pi }\left[
e^{F\left( h,\mathbf{x}\right) }\right] +\ln \frac{1}{\delta } \\
&&+\left( M+\beta m\right) \epsilon _{\beta }+\sum_{k=1}^{K}\left( \beta
_{k}-\beta _{k-1}\right) m\epsilon _{\beta _{k-1}}.
\end{eqnarray*}
\end{theorem}

The proof, given in Section \ref{Proofs stability}, is
similar to that of Theorem \ref{Theorem Main} and applies Markov's
inequality with $\nu _{\beta }$ instead of $G_{\beta }$. It then uses the
fact that, if $f$ is a bounded measurable function, then $\left\vert \mathbb{%
E}_{\nu _{1}}\left[ f\right] -\mathbb{E}_{\nu _{2}}\left[ f\right]
\right\vert \leq \left\Vert f\right\Vert _{\infty }d_{TV}\left( \nu _{1},\nu
_{2}\right) $. For the single-draw version (i)\ this is applied to $\left( 
\mathbb{E}_{\nu _{\beta }}-\mathbb{E}_{G_{\beta }}\right) \left[ e^{F\left(
h,\mathbf{x}\right) +\beta \hat{L}\left( h,\mathbf{x}\right) +\ln Z_{\beta
}\left( \mathbf{x}\right) }\right] $, which causes the exponential
dependence on $\beta m$ and $M$. For (ii) we can apply this in the exponent to 
$\left( \mathbb{E}_{\nu _{\beta }}-\mathbb{E}_{G_{\beta }}\right) \big[
F( h,\mathbf{x}) +\beta \hat{L}( h,\mathbf{x}) \big],$ and the logarithm makes the dependence linear. The rest of the proof is
mechanical.

The last terms in (i) and (ii) are the additional errors due to the
approximations of the Gibbs posteriors. The worst term is clearly the first
one in (i) due to the exponential dependence on the proxy function $F$,
which is typically of order $n$ and on $\beta $, which is larger than $n$ in
the regime in which we are interested. In practice, the requirement of such
an approximation is prohibitive for the single-draw version. The bound in
(ii)\ has more moderate approximation requirements. In this case, we can also
give bounds in terms of the $W_{2}$-Wasserstein metric (as
guaranteed by Corollary \ref{Corollary approximations}), if $F$ sand $\ell$ satisfy a Lipschitz
condition instead of boundedness.  We will use the
following fact: Since $W_{1}\leq W_{2}$ it follows from the
Kantorovich-Rubinstein Theorem \citep{villani2009optimal}, that for any
real Lipschitz function $f$ on $\mathcal{H}$ and probability measures $\nu
_{1},\nu _{2}\in \mathcal{P}\left( \mathcal{H}\right) $%
\begin{equation*}
\left\vert \mathbb{E}_{h\sim \nu _{1}}\left[ f\left( h\right) \right] -%
\mathbb{E}_{h\sim \nu _{2}}\left[ f\left( h\right) \right] \right\vert \leq
\left\Vert f\right\Vert _{\rm Lip}W_{1}\left( \nu _{1},\nu _{2}\right) \leq
\left\Vert f\right\Vert _{\rm Lip}W_{2}\left( \nu _{1},\nu _{2}\right) ,
\end{equation*}%
where $\left\Vert .\right\Vert _{\rm Lip}$ is the Lipschitz-seminorm. The
following result is then immediate, with proof exactly as in (ii) of Theorem %
\ref{Theorem BV approximations}.

\begin{theorem}
\label{Theorem Wasserstein}Assume the conditions of Theorem \ref{Theorem BV
approximations}, except that instead of $\left\vert \ell \left( h,\mathbf{x}%
\right) \right\vert \leq m$ and $\left\vert F\left( h,\mathbf{x}\right)
\right\vert \leq M$ we have $\left\Vert \ell \left( .,\mathbf{x}\right)
\right\Vert _{\rm Lip}\leq m$ and $\left\Vert F\left( .,\mathbf{x}\right)
\right\Vert _{\rm Lip}\leq M$ and that $W_{2}\left( \nu _{\beta _{k}}\left( 
\mathbf{x}\right) ,G_{\beta }\left( \mathbf{x}\right) \right) =\epsilon $.
Then, with probability at least $1-\delta $ as $\mathbf{x}\sim \mu ^{n}$ 
\begin{eqnarray*}
\mathbb{E}_{h\sim \nu _{\beta }\left( \mathbf{x}\right) }\left[ F\left( h,%
\mathbf{x}\right) \right] &\leq &\mathbb{E}_{h\sim \nu _{\beta }\left( 
\mathbf{x}\right) }\left[ \Gamma _{\nu }\left( h,\mathbf{x},\bm{\beta}%
\right) \right] +\ln \mathbb{E}_{\mathbf{x}}\mathbb{E}_{h\sim \pi }\left[
e^{F\left( h,\mathbf{x}\right) }\right] +\ln \frac{1}{\delta } \\
&&+\left( M+\beta m\right) \epsilon _{\beta }+\sum_{k=1}^{K}\left( \beta
_{k}-\beta _{k-1}\right) m\epsilon _{\beta _{k-1}}.
\end{eqnarray*}
\end{theorem}
\section{Experiments}
The purpose of our experiments is to show that our method gives nontrivial bounds on the test error for real-world data, while correctly bounding the test error on impossible data, where the same algorithm also achieves a small training error. For this demonstration, we use binary classification tasks. The real-world data are either the MNIST dataset, subdivided into the two classes of characters 0-4 and 5-9, or the CIFAR-10 dataset to distinguish between animals and vehicles. For impossible data, we randomize the labels of the training data. Our experiments are computationally heavy, so we generally use small sample sizes, from 2000 to 8000 examples. The hypothesis space is the set of weight vectors for a neural network with ReLU activation functions constrained by a Gaussian prior distribution with $\sigma =5$. Neural network architectures are described in Section \ref{sec:net_arch} of the appendix. To approximately sample the weight vectors in the vicinity of the Gibbs posterior, we use ULA as in (\ref{Our LMC}) or SGLD \citep{welling2011bayesian} with constant step size $\eta $. To ensure reproducibility, we provide the code and experimental results in an anonymous repository at \url{https://anonymous.4open.science/r/Gibbs-Generalization-45F1}.

\subsection{The loss function $\ell $}

Unless $\sigma$ is very small, which would prevent near-zero training
errors for random labels, an unbounded loss function like binary cross-entropy
entropy (BCE) would cause $\mathbb{E}_{h\sim G_{0}}\big[ \hat{L}(
h)\big] =\mathbb{E}_{h\sim \pi}\big[ \hat{L}(h)\big] $ and therefore the first term in the summation of $\Gamma$ in \eqref{eq:gamma} to be very large. For this reason we use a bounded loss function $\ell $, either bounded binary
cross-entropy as described in Appendix D of \cite{dziugaite2018data} or
the savage loss \citep{masnadi2008design}. We compute bounds 
%both for $\ell $ and the practically more relevant 
for the 0-1 loss, using the method described in Section \ref{Section
secondary loss}.

\subsection{Approximating the ergodic mean} \label{sec:ergodic_mean}
As we know of no sufficient criterion for convergence, we terminate
iterations at time $T$, when a very slow running mean $\mathbb{M}_{\rm stop}$ of
the loss trajectory $\big( \hat{L}( h_{\beta _{k}t},\mathbf{x})
\big) _{t=0}^{T}$ stops decreasing. A second running mean 
$\mathbb{M}_{\rm erg}$ is used as an approximation of the ergodic mean and thus of
expectations in the invariant distribution. We thus replace all expectations 
$\mathbb{E}_{h\sim G_{\beta _{k}}}\big[ \hat{L}( h,\mathbf{x})\big] $ occurring in the bounds by $\mathbb{M}_{\rm erg}\big[ \big( \hat{L}( h_{\beta _{k}t},\mathbf{x}) \big) _{t=0}^{T}\big] $. Both
running means $\mathbb{M}_{\rm stop}$ and $\mathbb{M}_{\rm erg}$ are implemented as
first-order, recursive lowpass filters described in Section \ref{sec:mov_avg} of the appendix.

\subsection{Calibration\label{Section calibration}} 
The theoretical bounds of Corollary \ref{Corollary approximations} in
combination with Theorem \ref{Theorem BV approximations} can only serve as guidance for the computation of practical bounds, because the quantities $R$ and $\alpha $ are impossible to
estimate in practice. But even if we assume these to be in the order of
unity, the bounds are too coarse to distinguish between different
temperatures with realistic stepsizes.\newline
A simple calculation shows that $KL\left( G_{\beta },G_{2\beta }\right) \leq
\beta \left( \mathbb{E}_{G_{\beta }}\left[ \hat{L}\right] -\mathbb{E}%
_{G_{2\beta }}\left[ \hat{L}\right] \right) \leq \beta \mathbb{E}_{G_{\beta
}}\left[ \hat{L}\right] $ (Lemma \ref{Lemma KLG2G} in Appendix \ref{Section
miscellaneous lemmata}). By Corollary \ref{Corollary approximations} we
should therefore have at least $8\eta dR^{2}/\alpha <\mathbb{E}_{G_{\beta
}\left( \mathbf{x}\right) }\left[ \hat{L}\right] $ to distinguish between
the expectations in the Gibbs posterior for $\beta $ and $2\beta $. The
smallest neural network we use has $d=392,500$. If $\ell $ has values in $%
\left[ 0,1\right] $ then $\mathbb{E}_{G_{\beta }\left( \mathbf{x}\right) }%
\left[ \hat{L}\right] \leq 1$, so even if $R$ and $\alpha $ are set to $1$,
we would need stepsizes in the order of $10^{-7}$. Safe values of $\eta $,
as suggested by the theoretical results in Section \ref{Langevin Theory},
are therefore impossible in practice, and the bound has to be adapted to a
realistic choice of $\eta $. 

The first possibility is to simply ignore the un-estimable error terms in
Theorem \ref{Theorem BV approximations}, and to compute the bounds from $%
\Gamma _{\nu }\left( h,\mathbf{x},\mathbf{\beta }\right) $ in place of $%
\Gamma \left( h,\mathbf{x},\mathbf{\beta }\right) $ as in Theorem \ref%
{Theorem Main}. We show some of these uncalibrated bounds, which already
make nontrivial predictions, in Section \ref{section uncalibrated}.

A second method uses a single
calibration parameter, whose value is computed from the training data. We
assume that the computed functional $\Gamma _{\nu }\left( h,\mathbf{x},%
\mathbf{\beta }\right) $ fails to estimate $\Gamma \left( h,\mathbf{x},%
\mathbf{\beta }\right) $ by a factor $r\left( \mathbf{x}\right) >0$, which
we compute as 
\begin{equation*}
r\left( \mathbf{x}\right) =\min \left\{ r:\forall k\in \left[ K\right] ,%
\text{ }\kappa ^{-1}\left( \mathbb{E}_{h\sim \nu _{\beta _{k}}\left( \mathbf{%
x}\right) }\left[ \hat{L}_{01}\left( h,\mathbf{\tilde{x}}\right) \right] ,%
\frac{1}{n}\left( r\Gamma _{\nu }\left( h,\mathbf{\tilde{x}},\mathbf{\beta }%
_{1}^{k}\right) +\ln \frac{2\sqrt{n}}{\delta }\right) \right) \geq \frac{1}{2%
}\right\} ,
\end{equation*}%
where $\mathbf{\tilde{x}}$ is the training set $\mathbf{x}$ with random
labels and $\hat{L}_{01}$ the empirical 01-error. The calibration value $r$
is thus the smallest factor of $\Gamma _{\nu }$, for which we obtain a
correct upper bound on the 01-error with random labels for all the $\beta
_{k}$.
By definition of $\Gamma
_{\nu }$, replacing it by $r \Gamma _{\nu }$ is equivalent to a multiplicative shift of the temperature scale.  \newline
It is a purely experimental finding, that our choice of $r$ leads to correct and surprisingly tight upper bounds on the test error of correctly labeled data in all cases we tried. We emphasize that our calibration procedure
depends only on the training data.

\subsection{Results} \label{sec:results}
\begin{figure}[htbp]
  \centering
  \begin{subfigure}[b]{0.425\linewidth} 
    \includegraphics[width=\linewidth]{figs/MRL2W1000SGLD8kLR001BBCE_01.png}
  \end{subfigure}  
  \begin{subfigure}[b]{0.425\linewidth} 
	\includegraphics[width=\linewidth]{figs/MCL2W1000SGLD8kLR001BBCE_01.png}
  \end{subfigure}  
    \begin{subfigure}[b]{0.425\linewidth} 
    \includegraphics[width=\linewidth]{figs/CRL2W1500SGLD8kLR0005BBCE_01.png}
  \end{subfigure}  
  \begin{subfigure}[b]{0.425\linewidth} 
	\includegraphics[width=\linewidth]{figs/CCL2W1500SGLD8kLR0005BBCE_01.png}
  \end{subfigure}  
  \caption{SGLD on MNIST and CIFAR-10 with 8000 training examples, MNIST above and CIFAR-10 below, random labels on the left, correct labels on the right. Both random and true labels are trained with the same algorithm and parameters on a fully connected ReLU network with two hidden layers of 1000 and 1500 units, respectively. The calibration factor for MNIST is 0.19, for CIFAR-10 0.22. Train error, test error and our bound for the Gibbs posterior average of the 0-1 loss are plotted against $\beta$.}
  \label{MC8k}
\end{figure}
Several experiments confirm the validity of the proposed bounds. An example is shown in Figure \ref{MC8k}, where a fully connected ReLU-network with two hidden layers of 1000 (respectively 1500) units each is trained with SGLD at inverse temperatures $\beta=$ 0, 500, 1000, 2000, 4000, 8000, 16000, 32000, and 64000. The train error for random labels is about 0.1 (or even less) at $\beta=64000$, where the bound is above 0.5. The test error for correct labels, however, is tightly bounded above. \\
Notice that for MNIST, which has the tightest bounds, the training error for the true labels is rapidly decreasing from 0.5 to 0.17 at $\beta=500$ and to 0.1 at $\beta=1000$.
The more moderate initial decrease for CIFAR-10 corresponds to the tendency to overfit on this more difficult dataset. This confirms the intuition, that good generalization at low temperatures is already announced in the high temperature regime. \\
We generally found the uncalibrated bounds for ULA tighter than those for SGLD, consistent with the findings of \cite{brosse2018promises}.
Experimental bounds for single draws from the posterior and 
various other experiments are reported  in Section \ref{sec:exp_res}. 

\section{Conclusion}

Using the integral representation of the log-partition function, the
Gibbs posterior admits the computation of upper bounds on the true error based on the training data and for any temperature.  These bounds are stable
under perturbation in the total-variation and Wasserstein metrics, and can
be approximated by Langevin Monte Carlo (LMC) algorithms. However, for realistic experiments, the approximations obtained by these algorithms
are coarse and require calibration, which leads to rather tight bounds in the interpolation regime of overparametrized neural networks. \\
The fact that the calibrated bounds are very tight is, at this point, a purely experimental finding, requiring more theoretical investigation in future work.

% \subsubsection*{Acknowledgments}
% Use unnumbered third-level headings for the acknowledgments. All
% acknowledgments, including those to funding agencies, go at the end of the paper.

%\massi{Many words, such as Langevin, in the references should be capitalized}
\bibliography{Bibfile}
\bibliographystyle{iclr2026_conference}

\appendix
\section*{Appendix}

In this appendix, we summarize a glossary of notation, give additional theoretical results and missing proofs, and provide more information on the numerical experiments, as well as additional experimental results.

\section{Table of notation}
\begin{tabular}{l|l|l}
\hline
Notation & Brief description & Section \\ \hline
$\mathcal{X}$ & space of data & \ref{Section preliminaries} \\ \hline
$\Sigma $ & sigma algebra (events) on $\mathcal{X}$ & \ref{Section
preliminaries} \\ \hline
$\mu $ & probability of data & \ref{Section preliminaries} \\ \hline
$n$ & sample size & \ref{Section introduction}, \ref{Section preliminaries} \\ \hline
$\mathbf{x}$ & generic member $\left( x_{1},...,x_{n}\right) \in \mathcal{X}%
^{n}$, training sample & \ref{Section preliminaries} \\ \hline
$\mathcal{H}$ & hypothesis space & \ref{Section preliminaries} \\ \hline
$\Omega $ & sigma algebra (events) on $\mathcal{H}$ & \ref{Section
preliminaries} \\ \hline
$\ell $ & $\ell :\mathcal{H\times X\rightarrow }\left[ 0,\infty \right) $
loss function & \ref{Section preliminaries} \\ \hline
$f$ & secondary loss function & \ref{Section
secondary loss} \\ \hline
$\mathcal{P}\left( \mathcal{H}\right) $ & probability measures on $\mathcal{H%
}$ & \ref{Section preliminaries} \\ \hline
$\pi $ & nonnegative a-priori measure on $\mathcal{H}$ & \ref{Section
preliminaries} \\ \hline
$\sigma$ & width of Gaussian prior & \ref{Langevin Theory}
\\ \hline $L\left( h\right) $ & $L\left( h\right) =\mathbb{E}_{x\sim \mu }\left[
h\left( x\right) \right] $, expected loss of $h\in \mathcal{H}$ & \ref%
{Section preliminaries} \\ \hline
$\hat{L}\left( h,\mathbf{x}\right) $ & $\hat{L}\left( h,\mathbf{x}\right)
=\left( 1/n\right) \sum_{i=1}^{n}\ell \left( h,x_{i}\right) $, empirical
loss of $h\in \mathcal{H}$ & \ref{Section preliminaries} \\ \hline
$\beta $ & inverse temperature & \ref{Section introduction}, \ref{Section preliminaries} \\ \hline
$Z_{\beta }\left( \mathbf{x}\right) $ & partition function & \ref{Section
preliminaries} \\ \hline
$G_{\beta ,\pi }\left( \mathbf{x}\right) $ & Gibbs posterior with energy $%
\hat{L}$ and prior $\pi $ & \ref{Section preliminaries} \\ \hline
$\mathbb{E}_{g\sim G_{\beta }\left( \mathbf{x}\right) }$ & posterior
expectation & \ref{Section preliminaries} \\ \hline
$\bm{\beta }$ & increasing sequence $\left( \beta _{1}<...<\beta
_{K}\right) $ of positive reals & \ref{Section integral representation} \\ 
\hline
$\Gamma \left( h,\mathbf{x},\bm{\beta}\right) $ & bounding functional & \ref{Section
integral representation} \\ \hline
$F\left( h,\mathbf{x}\right) $ & placeholder for generalization gap & \ref{Section bounds} \\ 
\hline
$\kappa $ & $kl\left( p,q\right) =p\ln \frac{p}{q}+\left( 1-p\right) \ln 
\frac{1-p}{1-q}$, rel. entropy of Bernoulli variables & \ref{Section
preliminaries} \\ \hline
$KL\left( \rho ,\nu \right) $ & $\int \left( \ln \frac{d\rho }{d\nu }\right)
d\rho $, KL-divergence of $\rho ,\nu \in \mathcal{P}\left( \mathcal{H}%
\right) $ & \ref{Section
preliminaries}, \ref{Section ULA},  \ref{Section stability} \\ \hline
$d_{TV}\left( \rho ,\nu \right) $ & total variation distance & \ref{Section
preliminaries}, \ref{Section ULA},  \ref{Section stability} \\ \hline
$W_{p}\left( \rho ,\nu \right) $ & $p$-Wasserstein distance & \ref{Section
preliminaries}, \ref{Section ULA},  \ref{Section stability} \\ \hline
$\Gamma _{\nu }\left( h,\mathbf{x},\bm{\beta}\right) $ & LMC
approximation of $\Gamma \left( h,\mathbf{x},\bm{\beta}\right) $ &  \ref{Section stability}\\ 
\hline
$\eta$ & step size or learning rate & \ref{Section ULA} \\ \hline
$\nu _{\beta ,\eta }$ & invariant measure of LMC approximation of $G_{\beta }
$ with step size $\eta $ & \ref{Section ULA} \\ \hline
$\nu _{\beta ,\eta ,t}$ & LMC approximation of $G_{\beta }$ with step size $%
\eta $ at iteration $t$ & \ref{Section ULA} \\ \hline
$r({\mathbf x})$ & calibration factor & \ref{Section calibration} \\ \hline
$\mathbf{\tilde{x}}$ & randomly labeled  data & \ref{Section calibration} \\ \hline
$\mathbb{M}_{\rm stop}$,$\mathbb{M}_{\rm erg}$ & filters for stopping and ergodic mean & \ref{sec:ergodic_mean}, \ref{sec:mov_avg} \\ \hline
\end{tabular}%
\section{Additional results and proofs}

\subsection{Sub-gaussian losses\label{Section Subgaussian}}

The freedom in the choice of $F$ allows a number of bounds to be derived
from Theorem \ref{Theorem Main}. A centered real random variable $Y$ is called $\sigma $-sub-Gaussian
if $\ln \mathbb{E}e^{\lambda Y-\mathbb{E}Y}\leq \lambda ^{2}\sigma ^{2}/2$
for all $\lambda \in \mathbb{R}$. Now, suppose that for some real function $f$
all the $x\in \mathcal{X\mapsto }f\left( h,x\right) $ are $\sigma $%
-sub-Gaussian as $x\sim \mu $. Let  $\hat{L}(h, \mathbf{x}) = \frac{1}{n}\sum_{i=1}^{n} f(h, x_i)$ and $L(h) = E_x[f(h, \mathbf{x})]$. Then, $\mathbf{x}\in \mathcal{X}^{n}\mapsto 
\hat{L}\left( h,\mathbf{x}\right) $ as $\mathbf{x}\sim \mu ^{n}$ is $\sigma /%
\sqrt{n}$-sub-Gaussian. It is tempting to set $F\left( h,\mathbf{x}\right)
=\lambda \left( L\left( h\right) -\hat{L}\left( h,\mathbf{x}\right) \right) $
in Theorem \ref{Theorem Main}, divide by $\lambda $ and then optimize over $%
\lambda $. Unfortunately, the last step is impossible, since the optimal $%
\lambda $ is data-dependent in its dependence on $\Gamma $ and ruins the
exponential moment bound on $F$. A more careful argument establishes the
following.

\begin{corollary}
	\label{Corollary Stratify}Suppose that for all $h\in \mathcal{H}$ the random
	variables $x\in \mathcal{X\mapsto} f \left( h,x\right) $ as $x\sim \mu $
	are $\sigma $-sub-Gaussian. For $\delta >0$ with probability at least $%
	1-\delta $ as $\mathbf{x}\sim \mu ^{n}$ and $h\sim G_{\beta ,\pi }\left( 
	\mathbf{x}\right) $ if $\Gamma \left( h,\mathbf{x},\bm{\beta}\right)
	\geq 1$, then,%
	\begin{equation*}
		E_{x}\left[ f\left( h,x\right) \right] -\frac{1}{n}\sum_{i}f\left(
		h,x_{i}\right) \leq \sigma \sqrt{\frac{2[\Gamma \left( h,\mathbf{x},\mathbf{%
					\beta }\right) \left( 1+1/n\right) +\ln \left( \Gamma \left( h,\mathbf{x},%
				\bm{\beta}\right) \left( n+1\right) /\delta \right)]}{n}}.
	\end{equation*}
\end{corollary}

For the proof, we use the following auxiliary result.

\begin{lemma}
	\label{Lemma model selection} \cite[][Lemma 15.6]{Anthony99} Suppose $\Pr $
	is a probability distribution and%
	\begin{equation*}
		\left\{ E\left( \alpha _{1},\alpha _{2},\delta \right) :0<\alpha _{1},\alpha
		_{2},\delta \leq 1\right\}
	\end{equation*}%
	is a set of events, such that
	
	(i) For all $0<\alpha \leq 1$ and $0<\delta \leq 1$,%
	\begin{equation*}
		\Pr \left\{ E\left( \alpha ,\alpha ,\delta \right) \right\} \leq \delta .
	\end{equation*}
	
	(ii) For all $0<\alpha _{1}\leq \alpha \leq \alpha _{2}\leq 1$ and $0<\delta
	_{1}\leq \delta \leq 1$%
	\begin{equation*}
		E\left( \alpha _{1},\alpha _{2},\delta _{1}\right) \subseteq E\left( \alpha
		,\alpha ,\delta \right) .
	\end{equation*}%
	Then for $0<a,\delta <1$,%
	\begin{equation*}
		\Pr \bigcup_{\alpha \in \left( 0,1\right] }E\left( \alpha a,\alpha ,\delta
		\alpha \left( 1-a\right) \right) \leq \delta .
	\end{equation*}
\end{lemma}

We put this lemma in a more convenient form.

\begin{lemma}
	\label{Lemma IvanovTychonov}Let $Y$ and $X\geq 0$ be real random variables, $%
	\psi :\mathbb{R}\times \left( 0,1\right) \rightarrow \mathbb{R}$ be
	increasing in the 1st argument and $\forall C>1$, $\delta \in \left(
	0,1\right) $, 
	\begin{equation*}
		\Pr \left\{ X\leq C\wedge Y>\psi \left( C,\delta \right) \right\} <\delta .
	\end{equation*}%
	Then, for every $\epsilon >0$%
	\begin{equation*}
		\Pr \left\{ X\geq 1\wedge Y>\psi \left( X\left( 1+\epsilon \right) ,\frac{%
			\delta \epsilon }{X\left( 1+\epsilon \right) }\right) \right\} .
	\end{equation*}
\end{lemma}

\begin{proof}
	This follows from Lemma \ref{Lemma model selection} using the events 
	\begin{equation*}
		E\left( \alpha _{1},\alpha _{2},\delta \right) =\left\{ X\leq \alpha
		_{2}^{-1}\wedge Y>f\left( \alpha _{1}^{-1},\delta \right) \right\} 
	\end{equation*}%
	and $a=1/\left( 1+\epsilon \right) $.
\end{proof}

\begin{proof}[Proof of Corollary \protect\ref{Corollary Stratify}]
	Take $F=\lambda \left( L\left( h\right) -\hat{L}\left( h,\mathbf{x}\right)
	\right) $, so $\ln \mathbb{E}_{\mathbf{x}}\mathbb{E}_{g\sim \pi }\left[
	e^{F\left( g,\mathbf{x}\right) }\right] \leq \lambda ^{2}\sigma ^{2}/\left(
	2n\right) $. From Theorem \ref{Theorem Main} and the properties of
	sub-Gaussian variables we get with $\lambda =\sigma ^{-1}\sqrt{2n\left(
		C+\ln \left( 1/\delta \right) \right) }$ that 
	\begin{eqnarray*}
		&&\Pr \left\{ \Gamma \left( h,\mathbf{x},\bm{\beta}\right) \leq C\wedge
		L\left( h\right) -\hat{L}\left( h,\mathbf{x}\right) >\sigma \sqrt{\frac{%
				2\left( C+\ln \left( 1/\delta \right) \right) }{n}}\right\}  \\
		&=&\Pr \left\{ \Gamma \left( h,\mathbf{x},\bm{\beta}\right) \leq
		C\wedge L\left( h\right) -\hat{L}\left( h,\mathbf{x}\right) >\frac{C+\ln
			\left( 1/\delta \right) }{\lambda }+\frac{\lambda \sigma ^{2}}{2n}\right\} 
		\\
		&=&\Pr \left\{ \Gamma \left( h,\mathbf{x},\bm{\beta}\right) \leq
		C\wedge \lambda \left( L\left( h\right) -\hat{L}\left( h,\mathbf{x}\right)
		\right) >C+\lambda ^{2}\sigma ^{2}/\left( 2n\right) +\ln \left( 1/\delta
		\right) \right\} \leq \delta .
	\end{eqnarray*}%
	Substitution in Lemma \ref{Lemma IvanovTychonov} with $\psi \left( C,\delta
	\right) =\sigma \sqrt{2\left( C+\ln \left( 1/\delta \right) \right) /n}$, $%
	X=\Gamma \left( h,\mathbf{x},\bm{\beta}\right) $, $Y=L\left( h\right) -%
	\hat{L}\left( h,\mathbf{x}\right) $ and $\epsilon =1/n$ gives Corollary \ref%
	{Corollary Stratify}.
\end{proof}

\bigskip 

\subsection{Proofs for Section \protect\ref{Section stability}\label{Proofs
		stability}}

Restatement of Theorem \ref{Theorem BV approximations}:

\begin{theorem}
	Suppose that $\mathcal{H}=\mathbb{R}^{d}$ and that there are numbers $%
	m,M<\infty $ such that for every $\mathbf{x}$ in $\mathcal{X}^{n}$ and $h\in 
	\mathcal{H}$ we have $\left\vert \ell \left( h,\mathbf{x}\right) \right\vert
	\leq m$ and $\left\vert F\left( h,\mathbf{x}\right) \right\vert \leq M$. Let 
	$0=\beta _{0}<\beta _{1}<\cdots<\beta _{K}=\beta $ and $\nu _{\beta _{k}}\left( 
	\mathbf{x}\right) \in \mathcal{P}\left( \mathcal{H}\right) $ be such that $%
	d_{TV}\left( \nu _{\beta _{k}}\left( \mathbf{x}\right) ,G_{\beta_k }\left( 
	\mathbf{x}\right) \right) =\epsilon _{\beta _{k}}$. Then,
	
	(i) with probability at least $1-\delta $ as $\mathbf{x}\sim \mu ^{n}$ and $%
	h\sim \nu _{\beta }\left( \mathbf{x}\right) $%
	\begin{eqnarray*}
		F\left( h,\mathbf{x}\right) &\leq &\Gamma _{\nu }\left( h,\mathbf{x},\mathbf{%
			\beta }\right) +\ln \mathbb{E}_{\mathbf{x}}\mathbb{E}_{h\sim \pi }\left[
		e^{F\left( h,\mathbf{x}\right) }\right] +\ln \frac{1}{\delta } \\
		&&+\ln \left( 2e^{M+\beta m}\epsilon _{\beta }\right) +\sum_{k=1}^{K}\left(
		\beta _{k}-\beta _{k-1}\right) m\epsilon _{\beta _{k-1}}.
	\end{eqnarray*}
	
	(ii) with probability at least $1-\delta $ as $\mathbf{x}\sim \mu ^{n}$ 
	\begin{eqnarray*}
		\mathbb{E}_{h\sim \nu _{\beta }\left( \mathbf{x}\right) }\left[ F\left( h,%
		\mathbf{x}\right) \right]  &\leq &\mathbb{E}_{h\sim \nu _{\beta }\left( 
			\mathbf{x}\right) }\left[ \Gamma _{\nu }\left( h,\mathbf{x},\bm{\beta}%
		\right) \right] +\ln \mathbb{E}_{\mathbf{x}}\mathbb{E}_{h\sim \pi }\left[
		e^{F\left( h,\mathbf{x}\right) }\right] +\ln \frac{1}{\delta } \\
		&&+\left( M+\beta m\right) \epsilon _{\beta }+\sum_{k=1}^{K}\left( \beta
		_{k}-\beta _{k-1}\right) m\epsilon _{\beta _{k-1}}.
	\end{eqnarray*}%
	
    \end{theorem}

\begin{proof}
	By the bound on $\ell $ we have%
	\begin{equation}
		\Gamma \left( h,\mathbf{x},\bm{\beta}\right) \leq \Gamma _{\nu }\left(
		h,\mathbf{x},\bm{\beta}\right) +\sum_{k=1}^{K}\left( \beta _{k}-\beta
		_{k-1}\right) m\epsilon _{\beta _{k-1}}\text{.}  \label{GammaGammaNuBound}
	\end{equation}%
	(i) From Markov's inequality we have (in analogy to the proof of Theorem \ref%
	{Theorem Main}) with probability at least $1-\delta $ as $x\sim \mu ^{n}$
	and $h\sim \nu _{\beta }\left( \mathbf{x}\right) $, that%
	\begin{eqnarray*}
		&&F\left( h,\mathbf{x}\right) +\beta \hat{L}\left( h,\mathbf{x}\right) +\ln
		Z_{\beta }\left( \mathbf{x}\right)  \\
		&\leq &\ln \mathbb{E}_{\mathbf{x}}\mathbb{E}_{h\sim \nu _{\beta }\left( 
			\mathbf{x}\right) }\left[ e^{F\left( h,\mathbf{x}\right) +\beta \hat{L}%
			\left( h,\mathbf{x}\right) +\ln Z_{\beta }\left( \mathbf{x}\right) }\right]
		+\ln \left( 1/\delta \right)  \\
		&\leq &\ln \left( \mathbb{E}_{\mathbf{x}}\mathbb{E}_{h\sim G_{\beta }\left( 
			\mathbf{x}\right) }\left[ e^{F\left( h,\mathbf{x}\right) +\beta \hat{L}%
			\left( h,\mathbf{x}\right) +\ln Z_{\beta }\left( \mathbf{x}\right) }\right]
		+e^{M+\beta m}\epsilon _{\beta }\right) +\ln \left( 1/\delta \right)  \\
		&=&\ln \left( \mathbb{E}_{\mathbf{x}}\mathbb{E}_{h\sim \pi }\left[
		e^{F\left( h,\mathbf{x}\right) }\right] +e^{M+\beta m}\epsilon _{\beta
		}\right) +\ln \left( 1/\delta \right)  \\
		&\leq &\ln \mathbb{E}_{\mathbf{x}}\mathbb{E}_{h\sim \pi }\left[ e^{F\left( h,%
			\mathbf{x}\right) }\right] +\ln \left( 2e^{M+\beta m}\epsilon _{\beta
		}\right) +\ln \left( 1/\delta \right) .
	\end{eqnarray*}%
	In the second inequality we used $\ln Z_{\beta }\left( \mathbf{x}\right)
	\leq 0$ and in the last line we used for $a,b\geq 1$ that $\ln \left(
	a+b\right) \leq \ln \max \left\{ a,b\right\} +\ln 2\leq \ln a+\ln b+\ln
	2=\ln a+\ln 2b$. Subtract $\beta \hat{L}\left( h,\mathbf{x}\right) +\ln
	Z_{\beta }\left( \mathbf{x}\right) $, use (\ref{density bound}) and (\ref%
	{GammaGammaNuBound}).
	
	(ii) Again with Markov's inequality, with probability at least $1-\delta $ as 
	$x\sim \mu ^{n}$%
	\begin{eqnarray*}
		&&\mathbb{E}_{h\sim \nu _{\beta }\left( \mathbf{x}\right) }\left[ F\left( h,%
		\mathbf{x}\right) +\beta \hat{L}\left( h,\mathbf{x}\right) +\ln Z_{\beta
		}\left( \mathbf{x}\right) \right] \\
		&\leq &\ln \mathbb{E}_{\mathbf{x}\sim \mu ^{n}}\left[ e^{\mathbb{E}_{h\sim
				\nu _{\beta }\left( \mathbf{x}\right) }\left[ F\left( h,\mathbf{x}\right)
			+\beta \hat{L}\left( h,\mathbf{x}\right) +\ln Z_{\beta }\left( \mathbf{x}%
			\right) \right] }\right] \\
		&\leq &\ln \mathbb{E}_{\mathbf{x}\sim \mu ^{n}}\left[ e^{\mathbb{E}_{h\sim
				G_{\beta }\left( \mathbf{x}\right) }\left[ F\left( h,\mathbf{x}\right)
			+\beta \hat{L}\left( h,\mathbf{x}\right) +\ln Z_{\beta }\left( \mathbf{x}%
			\right) \right] +\left( M+\beta m\right) \epsilon _{\beta }}\right] .
	\end{eqnarray*}%
	In the second inequality we used the fact that $\ln Z_{\beta }\left( \mathbf{%
		x}\right) \leq 0$ and the bounds on $F$ and $\ell $. Then, Jensen's
	inequality bounds the last line as%
	\begin{eqnarray*}
		&&\ln \mathbb{E}_{\mathbf{x}\sim \mu ^{n}}\mathbb{E}_{h\sim G_{\beta }\left( 
			\mathbf{x}\right) }\left[ e^{F\left( h,\mathbf{x}\right) +\beta \hat{L}%
			\left( h,\mathbf{x}\right) +\ln Z_{\beta }\left( \mathbf{x}\right) }\right]
		e^{\left( M+\beta m\right) \epsilon _{\beta }} \\
		&=&\ln \mathbb{E}_{\mathbf{x}\sim \mu ^{n}}\mathbb{E}_{h\sim \pi }\left[
		e^{F\left( h,\mathbf{x}\right) }\right] +\left( M+\beta m\right) \epsilon
		_{\beta }.
	\end{eqnarray*}%
	Again subtract $\beta \hat{L}\left( h,\mathbf{x}\right) +\ln Z_{\beta
	}\left( \mathbf{x}\right) $, use (\ref{density bound}) and (\ref%
	{GammaGammaNuBound}).
\end{proof}

\subsection{An ambiguity of the temperature scale\label{Section ambiguity}}

A multiplicative shift of the temperature scale is necessary to remove an
ambiguity, which arises from different scalings of the gradient. This can be
understood from Theorem \ref{Theorem BV approximations} and the nature of the update rule in (\ref{Our LMC}). Let $%
U_{\hat{L},\beta ,\eta }$ be the updating increment 
\[
U_{\hat{L},\beta ,\eta }\left( h\right) =-\eta \nabla \hat{L}\left( h\right)
-\frac{\eta h}{\beta \sigma ^{2}}+\sqrt{\frac{2\eta }{\beta }}\xi ,
\]%
Suppose that for a given loss landscape $\hat{L}$ and sufficiently small $%
\eta $, we have for all $\beta $ that $d_{TV}\left( G_{\beta },\nu _{\hat{L}%
,\beta ,\eta }\right) <2\sqrt{\frac{\eta d}{\beta \alpha }}\left( \beta R+%
\frac{1}{\sigma ^{2}}\right) <\epsilon $. Let another loss landscape $\hat{L}%
^{\prime }$ have a scaled gradient $\nabla \hat{L}^{\prime }=r\nabla \hat{L}$
with $r<1$, for example because $\hat{L}^{\prime }\left( h,\mathbf{x}\right)
=\hat{L}\left( rh,\mathbf{x}\right) $. Then, 
\[
U_{\hat{L}^{\prime },\beta ,\eta }\left( h\right) =-r\eta \nabla \hat{L}%
\left( h\right) -\frac{r\eta h}{r\beta \sigma ^{2}}+\sqrt{\frac{2r\eta }{%
r\beta }}\xi =U_{\hat{L},r\beta ,r\eta }\left( h\right) .
\]%
Since $r\eta <\eta $, by the bound in Theorem \ref{Theorem BV approximations} we should also have $%
d_{TV}\left( G_{r\beta },\nu _{\hat{L}^{\prime },\beta ,\eta }\right)
<\epsilon $ for all $\beta $. So the same algorithm approximates $G_{\beta }$
for $\hat{L}$ and $G_{r\beta }$ for $\hat{L}^{\prime }$. To compensate this shift in temperature scale we would have to multiply the $\Gamma$-functional for $\hat{L}^{\prime }$ with $r$. The calibration method
in Section \ref{Section calibration} is a convenient way to remove this ambiguity. 

The above suggests that smaller gradients tend to increase the $\Gamma$-functional and make the bounds more conservative. This may be the reason, why the shift, which guarantees an upper bound on random labels, in all cases also yields an upper bound for the true labels, because it seems intuitive that the loss landscape for true labels should tend to have smaller gradients than the one for random labels. This conjecture merits further theoretical and experimental investigation.

\subsection{Miscellaneous lemmata \label{Section miscellaneous lemmata}%
	\protect\bigskip}

\begin{lemma}
	\label{Lemma KLG2G}For $0<\beta <\infty $%
	\begin{equation*}
		\max \left\{ KL\left( G_{\beta },G_{2\beta }\right) ,KL\left( G_{2\beta
		},G_{2}\right) \right\} \leq \beta \left( \mathbb{E}_{h\sim G_{\beta }}\left[
		\hat{L}\left( h\right) \right] -\mathbb{E}_{h\sim G_{2\beta }}\left[ \hat{L}%
		\left( h\right) \right] \right)
	\end{equation*}%

\end{lemma}

\begin{proof}
 Using Lemma \ref{Lemma Derivatives}	\begin{eqnarray*}
		KL\left( G_{\beta },G_{2\beta }\right) &=&\mathbb{E}_{h\sim G_{\beta }}\left[
		-\beta \hat{L}\left( h\right) -\ln Z_{\beta }+2\beta \hat{L}\left( h\right)
		+\ln Z_{2\beta }\right] \\
		&=&\mathbb{E}_{h\sim G_{\beta }}\left[ \beta \hat{L}\left( h\right) \right]
		-\int_{\beta }^{2\beta }\mathbb{E}_{h\sim G_{\gamma }}\left[ \hat{L}\left(
		h\right) \right] d\gamma \\
		&\leq &\beta \left( \mathbb{E}_{h\sim G_{\beta }}\left[ \hat{L}\left(
		h\right) \right] -\mathbb{E}_{h\sim G_{2\beta }}\left[ \hat{L}\left(
		h\right) \right] \right)
	\end{eqnarray*}%
	Similarly%
	\begin{eqnarray*}
		KL\left( G_{2\beta },G_{2}\right) &=&\mathbb{E}_{h\sim G_{2\beta }}\left[
		-2\beta \hat{L}\left( h\right) -\ln Z_{2\beta }+\beta \hat{L}\left( h\right)
		+\ln Z_{\beta }\right] \\
		&=&-\mathbb{E}_{h\sim G_{\beta }}\left[ \beta \hat{L}\left( h\right) \right]
		+\int_{\beta }^{2\beta }\mathbb{E}_{h\sim G_{\gamma }}\left[ \hat{L}\left(
		h\right) \right] d\gamma \\
		&\leq &\beta \left( \mathbb{E}_{h\sim G_{\beta }}\left[ \hat{L}\left(
		h\right) \right] -\mathbb{E}_{h\sim G_{2\beta }}\left[ \hat{L}\left(
		h\right) \right] \right)
	\end{eqnarray*}
\end{proof}

\section{Experimental details and additional results}

\subsection{Experimental details}
All the codes to reproduce the results are provided through this \url{https://anonymous.4open.science/r/Gibbs-Generalization-45F1}. For all the experiments we use the isotropic Gaussian prior with $\mu =0$ and $\sigma = 5$. This induces an L2-regularization term in the energy function that is stated in the proof of Corollary~\ref{Corollary approximations}.

We use either standard SGLD or ULA with a constant step size and without additional correction terms. When ULA has been used, we use a step size of 0.01 for both datasets. However with SGLD, we set the step size to 0.01 for MNIST and 0.005 for CIFAR-10. For both datasets, MNIST and CIFAR-10, we use neural networks with ReLU activation functions.

\subsubsection{Network architecture} \label{sec:net_arch}
The fully connected networks consist of one, two, three hidden layers, each containing a constant number of units. Besides that we are using LeNet-5 architecture for MNIST and VGG16 architecture for CIFAR-10 to achieve low test error. For loss function $\ell $, we are using either bounded binary
cross-entropy (BBCE) as described in Appendix D of \cite{dziugaite2018data} or
the savage loss \citep{masnadi2008design}.

The LeNet-5 network follows a systematic pattern of alternating convolutional and pooling layers, followed by fully connected layers \citep{lecun2002gradient}. It begins with an input layer that accepts \(32 \times 32\) grayscale images. Thus, we pad our images to fit. 
The first convolutional layer (C1) applies 6 filters of size \(5 \times 5\) to extract low-level features, followed by a \(2 \times 2\) average pooling layer (S2) for spatial downsampling. The second convolutional layer (C3) uses 16 filters of size \(5 \times 5\) to capture more complex feature combinations, followed again by a \(2 \times 2\) average pooling layer (S4). 
A third convolutional layer (C5) with 120 filters of size \(5 \times 5\) acts as a feature extractor, producing 120 feature maps, each of size \(1 \times 1\). The architecture concludes with two fully connected layers: F6 with 84 neurons and a final output layer with 10 neurons for the original digit classification task. However, for our binary classification task, we modify F6 to have 420 neurons and use a single-neuron output layer. Throughout the network, ReLU activation functions replace the original tanh activations, which improves gradient flow and training performance in modern implementations.

VGG-16 is a widely used deep convolutional neural network architecture known for its simplicity and strong performance in image classification tasks \citep{simonyan2014very}. 
The architecture follows a consistent design using only \(3 \times 3\) convolutional filters and \(2 \times 2\) max pooling operations throughout the network. In our implementation, VGG-16 is adapted to handle CIFAR-10's smaller \(32 \times 32\) RGB images.  The network consists of 13 convolutional layers organized into five blocks: 
the first two blocks contain two convolutional layers each with 64 and 128 filters, respectively, while the last three blocks contain three convolutional layers each with 256, 512, and 512 filters, respectively. 
Each block is followed by a \(2 \times 2\) max pooling layer for spatial downsampling. All convolutional layers employ \(3 \times 3\) kernels with padding to preserve spatial dimensions, and ReLU activation functions introduce non-linearity.The convolutional feature extractor is followed by a classifier head consisting of three fully connected layers: two hidden layers with 1024 neurons each, using ReLU activation, and a final output layer with 1 neurons for binary classification. We also removed dropout to ensure that SGLD minimizes the defined energy function without any additional terms.

For MNIST, the input is a 784-dimensional vector, and the output is a scalar since we perform binary classification between digits 0--4 and 5--9. For CIFAR-10, the input dimension is 3072, and the output is again scalar, corresponding to binary classification between vehicles and animals. For evaluating our models, we are using all 10,000 test examples for both datasets.

\subsubsection{Minibatches} 
When using SGLD, we adopt minibatches of size proportional to $\sqrt{n}$. Thus, for $n=2000$ the mini-batch size is 50, and for $n=8000$ it is 100.
\subsubsection{Moving average filters} \label{sec:mov_avg}
As we explained in Section \ref{sec:ergodic_mean}, we are using a running mean $\mathbb{M}(x_1, \cdots, x_t)$ of $\hat{L}(h_j, \mathbf{x})$ from $j = 1, \cdots, t$ both as a criterion to stop the experiment and an estimation for $\mathbb{E}_{h\sim G_{\beta _{k}}}\left[ \hat{L}\left( h,\mathbf{x}\right) %
\right] $. We define the running mean recursively in one of two ways:

\begin{align*}
\mathbb{M}_t &= \tfrac{\alpha}{2} \hat{L}(h_t, \mathbf{x}) 
             + \tfrac{\alpha}{2} \hat{L}(h_{t-1}, \mathbf{x}) 
             + (1-\alpha)\,\mathbb{M}_{t-1}, \\
\mathbb{M}_t &= \alpha \hat{L}(h_t, \mathbf{x}) 
             + (1-\alpha)\,\mathbb{M}_{t-1},
\end{align*}
with $\mathbb{M}_0 = 0$ and small $\alpha$. 
We use the first (symmetric) form in the experiments with ULA, and 
the second (standard exponential moving average) form with SGLD 
for convenience. We set different values of $\alpha$ for the two roles: 
$\alpha = 0.0025$ for the stopping criterion ($\mathbb{M}_{\mathrm{stop}}$) 
and $\alpha = 0.01$ for approximating the ergodic mean ($\mathbb{M}_{\mathrm{erg}}$). 
The stopping rule is triggered when
\[
\mathbb{M}_t - \mathbb{M}_{t-1} \geq \epsilon,
\]
with $\epsilon = 10^{-7}$. 
To avoid premature termination, we impose a minimum of 4000 steps 
before applying this criterion. 
As $\alpha \to 0$ and $t \to \infty$, the quantity 
$\mathbb{M}_t$ converges to the ergodic mean.
% The above update rule is equivalent to a causal, first-order discrete-time low-pass filter with unit DC gain. The transfer function is
% $$H(z) = \frac{Y(z)}{X(z)} = \frac{\alpha}{1 - (1-\alpha)\,z^{-1}}$$
% and its frequency response on the unit circle $z = e^{j\omega}$ is
% $$H(e^{j\omega}) = \frac{\alpha}{1 - (1-\alpha)e^{-j\omega}}.$$
% % A convenient form for the magnitude is

% % $|H(e^{j\omega})| = \frac{\alpha}{\sqrt{(1-(1-\alpha)\cos\omega)^{2} + ((1-\alpha)\sin\omega)^{2}}} = \frac{\alpha}{\sqrt{1+(1-\alpha)^{2}-2(1-\alpha)\cos\omega}}.$

\subsection{Experimental results} \label{sec:exp_res}
\subsubsection{Illustration of Bound Computation} \label{sec:illustration}
In this section, we demonstrate again the figure in the main body in more details. 
The figure \ref{MC8k_kl} illustrates how our bounds are computed. The sequence of mean
training losses in $\ell $ is used to compute for each $\beta $ the
functional $\Gamma $ and the "KL-Bound", which corresponds to the right hand
side of the inequalities in Corollary \ref{Corollary kl}. Our bound on the test loss is
then computed by applying the function $\kappa ^{-1}$ to the empirical 0-1
error and to this kl-bound. The graph of "KL(Train, Test)" corresponds to the
left hand side in Corollary \ref{Corollary kl}.
\newline
It is remarkable that the close fit of the upper bound on the random labels is achieved by the adjustment of a single calibration parameter.

\begin{figure}[htbp]
  \centering
  \begin{subfigure}[b]{0.425\linewidth} 
    \includegraphics[width=\linewidth]{figs/MRL2W1000SGLD8kLR001BBCE_01_KL.png}
  \end{subfigure}  
  \begin{subfigure}[b]{0.425\linewidth} 
	\includegraphics[width=\linewidth]{figs/MCL2W1000SGLD8kLR001BBCE_01_KL.png}
  \end{subfigure}  
    \begin{subfigure}[b]{0.425\linewidth} 
    \includegraphics[width=\linewidth]{figs/CRL2W1500SGLD8kLR0005BBCE_01_KL.png}
  \end{subfigure}  
  \begin{subfigure}[b]{0.425\linewidth} 
	\includegraphics[width=\linewidth]{figs/CCL2W1500SGLD8kLR0005BBCE_01_KL.png}
  \end{subfigure}  
  \caption{A more detailed version of Figure \ref{MC8k} to illustrate how the bounds are computed.}
  \label{MC8k_kl}
\end{figure}

\subsubsection{Single-draws}
For the setting described in Section~\ref{sec:results}, we also present the bounds for the single-draw case in Figure~\ref{MC8k_singledraw}. It is noteworthy that, although the theoretical guarantees for this scenario are rather weak, the empirical bounds behave well. However, as visible in the plots, the results exhibit fluctuations and irregularities caused by stochastic effects, which make them less reliable.
\begin{figure}[htbp]
  \centering
  \begin{subfigure}[b]{0.425\linewidth} 
    \includegraphics[width=\linewidth]{figs/MRL2W1000SGLD8kLR001BBCE_01_singledraw.png}
  \end{subfigure}  
  \begin{subfigure}[b]{0.425\linewidth} 
	\includegraphics[width=\linewidth]{figs/MCL2W1000SGLD8kLR001BBCE_01_singledraw.png}
  \end{subfigure}  
    \begin{subfigure}[b]{0.425\linewidth} 
    \includegraphics[width=\linewidth]{figs/CRL2W1500SGLD8kLR0005BBCE_01_singledraw.png}
  \end{subfigure}  
  \begin{subfigure}[b]{0.425\linewidth} 
	\includegraphics[width=\linewidth]{figs/CCL2W1500SGLD8kLR0005BBCE_01_singledraw.png}
  \end{subfigure}  
  \caption{SGLD on MNIST and CIFAR-10 with 2000 training examples using BBCE loss function. The first row corresponds to MNIST and the second row to CIFAR-10. Random labels are shown on the left, correct labels on the right. Both random and true labels are trained with exactly the same algorithm and parameters on a fully connected ReLU network with two hidden layers of 1000 (respectively 1500) units. The calibration factor for MNIST is 0.77, for CIFAR-10 0.89. Train error, test error and our bound for a single-draw of the 0-1 loss are plotted against $\beta$.}
  \label{MC8k_singledraw}
\end{figure}
\FloatBarrier

\subsubsection{Different Architectures}
In this section, we evaluate the performance of different models and architectures on both MNIST and CIFAR-10, demonstrating that our bound can be used to guide model selection. In addition to the two-hidden-layer neural networks described in Section~\ref{MC8k}, we consider fully connected neural networks with three hidden layers, containing 500 and 1000 units for MNIST and CIFAR-10, respectively. Furthermore, we employ the LeNet-5 architecture for MNIST and VGG-16 for CIFAR-10 to achieve high test accuracy. Detailed descriptions of these architectures are provided in Section~\ref{sec:net_arch}.

Figure~\ref{MC8k_more} demonstrates the robustness of our bound across different models.  We observe that the bounds can be very tight even when the test error is small. For convolutional neural networks, especially on the MNIST dataset, we observe strong performance with the true labels, but relatively poor performance with random labels, despite having more parameters than training examples. This can be explained by the fact that convolutional architectures are still far from being highly overparameterized. For the MNIST dataset, we use fully connected neural networks with two or three hidden layers, containing 1000 or 500 units per layer, respectively. This corresponds to a total of approximately 1,787,000 and 893,000 parameters, resulting in a parameter-to-training-example ratio of roughly 200 and 100, respectively. In contrast, LeNet-5 has around 100,000 parameters, yielding a ratio of approximately 12.5.

\begin{figure}[htbp]
  \centering
  \begin{subfigure}[b]{0.425\linewidth} 
    \includegraphics[width=\linewidth]{figs/MRL3W500SGLD8kLR001BBCE_01.png}
  \end{subfigure}  
  \begin{subfigure}[b]{0.425\linewidth} 
	\includegraphics[width=\linewidth]{figs/MCL3W500SGLD8kLR001BBCE_01.png}
  \end{subfigure}  
    \begin{subfigure}[b]{0.425\linewidth} 
    \includegraphics[width=\linewidth]{figs/MRLLW500SGLD8kLR001BBCE_01.png}
  \end{subfigure}  
  \begin{subfigure}[b]{0.425\linewidth} 
	\includegraphics[width=\linewidth]{figs/MCLLW500SGLD8kLR001BBCE_01.png}
  \end{subfigure}  
    \begin{subfigure}[b]{0.425\linewidth} 
    \includegraphics[width=\linewidth]{figs/CRL3W1000SGLD8kLR0005BBCE_01.png}
  \end{subfigure}  
  \begin{subfigure}[b]{0.425\linewidth} 
	\includegraphics[width=\linewidth]{figs/CCL3W1000SGLD8kLR0005BBCE_01.png}
  \end{subfigure}  
    \begin{subfigure}[b]{0.425\linewidth} 
    \includegraphics[width=\linewidth]{figs/CRLVW500SGLD8kLR0005BBCE_01.png}
  \end{subfigure}  
  \begin{subfigure}[b]{0.425\linewidth} 
	\includegraphics[width=\linewidth]{figs/CCLVW500SGLD8kLR0005BBCE_01.png}
  \end{subfigure}  
  
  \caption{SGLD on MNIST and CIFAR-10 with 8000 training examples using BBCE loss function. The first two rows correspond to MNIST, and the remaining rows to CIFAR-10. Random labels are shown on the left, and correct labels on the right. Both random and true labels are trained using the same algorithm and hyperparameters on a fully connected ReLU network with three hidden layers of 500 (MNIST) or 1000 (CIFAR-10) units, followed by LeNet-5 (MNIST) or VGG-16 (CIFAR-10) shown in the subsequent row. The calibration factors for MNIST are 1.05 and 0.32, for CIFAR-10 0.96 and 0.73. The training error, test error, and our bound for the Gibbs posterior average of the 0–1 loss are plotted against \(\beta\).}
  \label{MC8k_more}
\end{figure}

The empirical test bounds can serve as a selection criterion among different models. Table \ref{tab:combined_results} show that test bounds at low temperature are useful for model selection, and that bounds at high temperature can also predict the behavior of the model at low temperature.


\begin{table}[htbp]
\centering

\begin{subtable}[t]{0.8\textwidth}
\centering
\begin{tabular}{lccc}
\toprule
 & 2HL (W=1000) & 3HL (W=500) & LeNet-5 \\
\midrule
Test Bound at $\beta = 1k$ & 0.1766 & 0.2347 & 0.0887 \\
Test Error at $\beta = 64k$ & 0.0498 & 0.0549 & 0.0317 \\
Test Bound at $\beta = 64k$ & 0.0860 & 0.1314 & 0.0375 \\
\bottomrule
\end{tabular}
\caption{MNIST, 8k training examples (true labels).}
\label{tab:mnist_results}
\end{subtable}

\vspace{0.5em}

\begin{subtable}[t]{0.8\textwidth}
\centering
\begin{tabular}{lccc}
\toprule
 & 2HL (W=1500) & 3HL (W=1000) & VGG-16 \\
\midrule
Test Bound at $\beta = 1k$ & 0.2905 & 0.3635 & 0.2330 \\
Test Error at $\beta = 64k$ & 0.1719 & 0.1782 & 0.0903 \\
Test Bound at $\beta = 64k$ & 0.2266 & 0.2807 & 0.2030 \\
\bottomrule
\end{tabular}
\caption{CIFAR-10, 8k training examples (true labels).}
\label{tab:cifar_results}
\end{subtable}

\caption{Test bounds and test errors for different neural network architectures on MNIST and CIFAR-10. The bounds at both low and high temperatures reliably reflect test error performance at low temperature.}
\label{tab:combined_results}
\end{table}
\FloatBarrier

\subsubsection{ULA}
We have also conducted experiments using ULA for both datasets. The main difference from SGLD is that we use all the information to compute the gradient at each step. The results are shown in Figure~\ref{MC2k_ULA_BBCE}. 

\begin{figure}[htbp]
  \centering
  \begin{subfigure}[b]{0.425\linewidth} 
    \includegraphics[width=\linewidth]{figs/MRL1W500ULA2kLR001BBCE_01.png}
  \end{subfigure}  
  \begin{subfigure}[b]{0.425\linewidth} 
	\includegraphics[width=\linewidth]{figs/MCL1W500ULA2kLR001BBCE_01.png}
  \end{subfigure}  
    \begin{subfigure}[b]{0.425\linewidth} 
    \includegraphics[width=\linewidth]{figs/CRL2W1000ULA2kLR001BBCE_01.png}
  \end{subfigure}  
  \begin{subfigure}[b]{0.425\linewidth} 
	\includegraphics[width=\linewidth]{figs/CCL2W1000ULA2kLR001BBCE_01.png}
  \end{subfigure}  
  \caption{ULA on MNIST and CIFAR-10 with 2000 training examples using BBCE loss function. The first row corresponds to MNIST and the second row to CIFAR-10. Random labels are shown on the left, correct labels on the right. Both random and true labels are trained with the same algorithm and parameters on a fully connected ReLU network with one (respectively two) hidden layers of 500 (respectively 1000) units. The calibration factor for MNIST is 1.98, for CIFAR-10 1.85. Train error, test error and our bound for the Gibbs posterior average of the 0-1 loss are plotted against $\beta$.}
  \label{MC2k_ULA_BBCE}
\end{figure}
\FloatBarrier

\subsubsection{Savage Loss function}
We additionally performed experiments using the Savage loss to verify the robustness of our results across different loss functions. Following the same setup as in the previous section, the outcomes are reported in Figure~\ref{MC2k_ULA_SAVAGE}. 
\begin{figure}[htbp]
  \centering
  \begin{subfigure}[b]{0.425\linewidth} 
    \includegraphics[width=\linewidth]{figs/MRL1W500ULA2kLR001SAVAGE_01.png}
  \end{subfigure}  
  \begin{subfigure}[b]{0.425\linewidth} 
	\includegraphics[width=\linewidth]{figs/MCL1W500ULA2kLR001SAVAGE_01.png}
  \end{subfigure}  
    \begin{subfigure}[b]{0.425\linewidth} 
    \includegraphics[width=\linewidth]{figs/CRL2W1000ULA2kLR001SAVAGE_01.png}
  \end{subfigure}  
  \begin{subfigure}[b]{0.425\linewidth} 
	\includegraphics[width=\linewidth]{figs/CCL2W1000ULA2kLR001SAVAGE_01.png}
  \end{subfigure}  
  \caption{ULA on MNIST and CIFAR-10 with 2000 training examples using savage loss function. The first row corresponds to MNIST and the second row to CIFAR-10. Random labels are shown on the left, correct labels on the right. Both random and true labels are trained with the same algorithm and parameters on a fully connected ReLU network with one (respectively two) hidden layers of 500 (respectively 1000) units. The calibration factor for MNIST is 0.49, for CIFAR-10 0.59. Train error, test error and our bound for the Gibbs posterior average of the 0-1 loss are plotted against $\beta$.}
  \label{MC2k_ULA_SAVAGE}
\end{figure}
\FloatBarrier
\subsubsection{Real-world use cases}

We further evaluated Stochastic Gradient Descent (SGD) to examine the practical relevance of our bounds in real-world interpolation regimes. 

Based on our observations, we suggest the following procedure for practitioners who wish to train overparameterized neural networks with standard SGD while also obtaining generalization guarantees. First, randomly permute the labels, train the network at different temperatures, and compute the bound together with the calibration factor. Then, repeat the same procedure using the true labels. At very low temperatures, this approach provides generalization guarantees that may transfer to SGD. The corresponding results are presented in Table~\ref{tab:combined_results_sgd}.

\begin{table}[htbp]
\centering

\begin{subtable}[t]{0.8\textwidth}
\centering
\begin{tabular}{lccc}
\toprule
 & 2HL (W=1000) & 3HL (W=500) & LeNet-5 \\
\midrule
Test Error with SGD & 0.0364 & 0.0363 & 0.0308 \\
Test Error at $\beta = 64k$ & 0.0498 & 0.0549 & 0.0317 \\
Test Bound at $\beta = 64k$ & 0.0860 & 0.1314 & 0.0375 \\
\bottomrule
\end{tabular}
\caption{MNIST, 8k training examples (true labels).}
\label{tab:mnist_results_sgd}
\end{subtable}

\vspace{0.5em}

\begin{subtable}[t]{0.8\textwidth}
\centering
\begin{tabular}{lccc}
\toprule
 & 2HL (W=1500) & 3HL (W=1000) & VGG-16 \\
\midrule
Test Error with SGD & 0.1423 & 0.1415 & 0.0933 \\
Test Error at $\beta = 64k$ & 0.1719 & 0.1782 & 0.0903 \\
Test Bound at $\beta = 64k$ & 0.2266 & 0.2807 & 0.2030 \\
\bottomrule
\end{tabular}
\caption{CIFAR-10, 8k training examples (true labels).}
\label{tab:cifar_results_sgd}
\end{subtable}

\caption{Comparing SGD test error with SGLD test errors and bounds for different neural network architectures on MNIST and CIFAR-10.}
\label{tab:combined_results_sgd}
\end{table}
\FloatBarrier
\end{document}
